\documentclass[11pt]{article}

    \usepackage[breakable]{tcolorbox}
    \usepackage{parskip} % Stop auto-indenting (to mimic markdown behaviour)
    

    % Basic figure setup, for now with no caption control since it's done
    % automatically by Pandoc (which extracts ![](path) syntax from Markdown).
    \usepackage{graphicx}
    % Keep aspect ratio if custom image width or height is specified
    \setkeys{Gin}{keepaspectratio}
    % Maintain compatibility with old templates. Remove in nbconvert 6.0
    \let\Oldincludegraphics\includegraphics
    % Ensure that by default, figures have no caption (until we provide a
    % proper Figure object with a Caption API and a way to capture that
    % in the conversion process - todo).
    \usepackage{caption}
    \DeclareCaptionFormat{nocaption}{}
    \captionsetup{format=nocaption,aboveskip=0pt,belowskip=0pt}

    \usepackage{float}
    \floatplacement{figure}{H} % forces figures to be placed at the correct location
    \usepackage{xcolor} % Allow colors to be defined
    \usepackage{enumerate} % Needed for markdown enumerations to work
    \usepackage{geometry} % Used to adjust the document margins
    \usepackage{amsmath} % Equations
    \usepackage{amssymb} % Equations
    \usepackage{textcomp} % defines textquotesingle
    % Hack from http://tex.stackexchange.com/a/47451/13684:
    \AtBeginDocument{%
        \def\PYZsq{\textquotesingle}% Upright quotes in Pygmentized code
    }
    \usepackage{upquote} % Upright quotes for verbatim code
    \usepackage{eurosym} % defines \euro

    \usepackage{iftex}
    \ifPDFTeX
        \usepackage[T1]{fontenc}
        \IfFileExists{alphabeta.sty}{
              \usepackage{alphabeta}
          }{
              \usepackage[mathletters]{ucs}
              \usepackage[utf8x]{inputenc}
          }
    \else
        \usepackage{fontspec}
        \usepackage{unicode-math}
    \fi

    \usepackage{fancyvrb} % verbatim replacement that allows latex
    \usepackage{grffile} % extends the file name processing of package graphics
                         % to support a larger range
    \makeatletter % fix for old versions of grffile with XeLaTeX
    \@ifpackagelater{grffile}{2019/11/01}
    {
      % Do nothing on new versions
    }
    {
      \def\Gread@@xetex#1{%
        \IfFileExists{"\Gin@base".bb}%
        {\Gread@eps{\Gin@base.bb}}%
        {\Gread@@xetex@aux#1}%
      }
    }
    \makeatother
    \usepackage[Export]{adjustbox} % Used to constrain images to a maximum size
    \adjustboxset{max size={0.9\linewidth}{0.9\paperheight}}

    % The hyperref package gives us a pdf with properly built
    % internal navigation ('pdf bookmarks' for the table of contents,
    % internal cross-reference links, web links for URLs, etc.)
    \usepackage{hyperref}
    % The default LaTeX title has an obnoxious amount of whitespace. By default,
    % titling removes some of it. It also provides customization options.
    \usepackage{titling}
    \usepackage{longtable} % longtable support required by pandoc >1.10
    \usepackage{booktabs}  % table support for pandoc > 1.12.2
    \usepackage{array}     % table support for pandoc >= 2.11.3
    \usepackage{calc}      % table minipage width calculation for pandoc >= 2.11.1
    \usepackage[inline]{enumitem} % IRkernel/repr support (it uses the enumerate* environment)
    \usepackage[normalem]{ulem} % ulem is needed to support strikethroughs (\sout)
                                % normalem makes italics be italics, not underlines
    \usepackage{soul}      % strikethrough (\st) support for pandoc >= 3.0.0
    \usepackage{mathrsfs}
    

    
    % Colors for the hyperref package
    \definecolor{urlcolor}{rgb}{0,.145,.698}
    \definecolor{linkcolor}{rgb}{.71,0.21,0.01}
    \definecolor{citecolor}{rgb}{.12,.54,.11}

    % ANSI colors
    \definecolor{ansi-black}{HTML}{3E424D}
    \definecolor{ansi-black-intense}{HTML}{282C36}
    \definecolor{ansi-red}{HTML}{E75C58}
    \definecolor{ansi-red-intense}{HTML}{B22B31}
    \definecolor{ansi-green}{HTML}{00A250}
    \definecolor{ansi-green-intense}{HTML}{007427}
    \definecolor{ansi-yellow}{HTML}{DDB62B}
    \definecolor{ansi-yellow-intense}{HTML}{B27D12}
    \definecolor{ansi-blue}{HTML}{208FFB}
    \definecolor{ansi-blue-intense}{HTML}{0065CA}
    \definecolor{ansi-magenta}{HTML}{D160C4}
    \definecolor{ansi-magenta-intense}{HTML}{A03196}
    \definecolor{ansi-cyan}{HTML}{60C6C8}
    \definecolor{ansi-cyan-intense}{HTML}{258F8F}
    \definecolor{ansi-white}{HTML}{C5C1B4}
    \definecolor{ansi-white-intense}{HTML}{A1A6B2}
    \definecolor{ansi-default-inverse-fg}{HTML}{FFFFFF}
    \definecolor{ansi-default-inverse-bg}{HTML}{000000}

    % common color for the border for error outputs.
    \definecolor{outerrorbackground}{HTML}{FFDFDF}

    % commands and environments needed by pandoc snippets
    % extracted from the output of `pandoc -s`
    \providecommand{\tightlist}{%
      \setlength{\itemsep}{0pt}\setlength{\parskip}{0pt}}
    \DefineVerbatimEnvironment{Highlighting}{Verbatim}{commandchars=\\\{\}}
    % Add ',fontsize=\small' for more characters per line
    \newenvironment{Shaded}{}{}
    \newcommand{\KeywordTok}[1]{\textcolor[rgb]{0.00,0.44,0.13}{\textbf{{#1}}}}
    \newcommand{\DataTypeTok}[1]{\textcolor[rgb]{0.56,0.13,0.00}{{#1}}}
    \newcommand{\DecValTok}[1]{\textcolor[rgb]{0.25,0.63,0.44}{{#1}}}
    \newcommand{\BaseNTok}[1]{\textcolor[rgb]{0.25,0.63,0.44}{{#1}}}
    \newcommand{\FloatTok}[1]{\textcolor[rgb]{0.25,0.63,0.44}{{#1}}}
    \newcommand{\CharTok}[1]{\textcolor[rgb]{0.25,0.44,0.63}{{#1}}}
    \newcommand{\StringTok}[1]{\textcolor[rgb]{0.25,0.44,0.63}{{#1}}}
    \newcommand{\CommentTok}[1]{\textcolor[rgb]{0.38,0.63,0.69}{\textit{{#1}}}}
    \newcommand{\OtherTok}[1]{\textcolor[rgb]{0.00,0.44,0.13}{{#1}}}
    \newcommand{\AlertTok}[1]{\textcolor[rgb]{1.00,0.00,0.00}{\textbf{{#1}}}}
    \newcommand{\FunctionTok}[1]{\textcolor[rgb]{0.02,0.16,0.49}{{#1}}}
    \newcommand{\RegionMarkerTok}[1]{{#1}}
    \newcommand{\ErrorTok}[1]{\textcolor[rgb]{1.00,0.00,0.00}{\textbf{{#1}}}}
    \newcommand{\NormalTok}[1]{{#1}}

    % Additional commands for more recent versions of Pandoc
    \newcommand{\ConstantTok}[1]{\textcolor[rgb]{0.53,0.00,0.00}{{#1}}}
    \newcommand{\SpecialCharTok}[1]{\textcolor[rgb]{0.25,0.44,0.63}{{#1}}}
    \newcommand{\VerbatimStringTok}[1]{\textcolor[rgb]{0.25,0.44,0.63}{{#1}}}
    \newcommand{\SpecialStringTok}[1]{\textcolor[rgb]{0.73,0.40,0.53}{{#1}}}
    \newcommand{\ImportTok}[1]{{#1}}
    \newcommand{\DocumentationTok}[1]{\textcolor[rgb]{0.73,0.13,0.13}{\textit{{#1}}}}
    \newcommand{\AnnotationTok}[1]{\textcolor[rgb]{0.38,0.63,0.69}{\textbf{\textit{{#1}}}}}
    \newcommand{\CommentVarTok}[1]{\textcolor[rgb]{0.38,0.63,0.69}{\textbf{\textit{{#1}}}}}
    \newcommand{\VariableTok}[1]{\textcolor[rgb]{0.10,0.09,0.49}{{#1}}}
    \newcommand{\ControlFlowTok}[1]{\textcolor[rgb]{0.00,0.44,0.13}{\textbf{{#1}}}}
    \newcommand{\OperatorTok}[1]{\textcolor[rgb]{0.40,0.40,0.40}{{#1}}}
    \newcommand{\BuiltInTok}[1]{{#1}}
    \newcommand{\ExtensionTok}[1]{{#1}}
    \newcommand{\PreprocessorTok}[1]{\textcolor[rgb]{0.74,0.48,0.00}{{#1}}}
    \newcommand{\AttributeTok}[1]{\textcolor[rgb]{0.49,0.56,0.16}{{#1}}}
    \newcommand{\InformationTok}[1]{\textcolor[rgb]{0.38,0.63,0.69}{\textbf{\textit{{#1}}}}}
    \newcommand{\WarningTok}[1]{\textcolor[rgb]{0.38,0.63,0.69}{\textbf{\textit{{#1}}}}}
    \makeatletter
    \newsavebox\pandoc@box
    \newcommand*\pandocbounded[1]{%
      \sbox\pandoc@box{#1}%
      % scaling factors for width and height
      \Gscale@div\@tempa\textheight{\dimexpr\ht\pandoc@box+\dp\pandoc@box\relax}%
      \Gscale@div\@tempb\linewidth{\wd\pandoc@box}%
      % select the smaller of both
      \ifdim\@tempb\p@<\@tempa\p@
        \let\@tempa\@tempb
      \fi
      % scaling accordingly (\@tempa < 1)
      \ifdim\@tempa\p@<\p@
        \scalebox{\@tempa}{\usebox\pandoc@box}%
      % scaling not needed, use as it is
      \else
        \usebox{\pandoc@box}%
      \fi
    }
    \makeatother

    % Define a nice break command that doesn't care if a line doesn't already
    % exist.
    \def\br{\hspace*{\fill} \\* }
    % Math Jax compatibility definitions
    \def\gt{>}
    \def\lt{<}
    \let\Oldtex\TeX
    \let\Oldlatex\LaTeX
    \renewcommand{\TeX}{\textrm{\Oldtex}}
    \renewcommand{\LaTeX}{\textrm{\Oldlatex}}
    % Document parameters
    % Document title
    \title{apresentacao\_orientadores}
    
    
    
    
    
    
    
% Pygments definitions
\makeatletter
\def\PY@reset{\let\PY@it=\relax \let\PY@bf=\relax%
    \let\PY@ul=\relax \let\PY@tc=\relax%
    \let\PY@bc=\relax \let\PY@ff=\relax}
\def\PY@tok#1{\csname PY@tok@#1\endcsname}
\def\PY@toks#1+{\ifx\relax#1\empty\else%
    \PY@tok{#1}\expandafter\PY@toks\fi}
\def\PY@do#1{\PY@bc{\PY@tc{\PY@ul{%
    \PY@it{\PY@bf{\PY@ff{#1}}}}}}}
\def\PY#1#2{\PY@reset\PY@toks#1+\relax+\PY@do{#2}}

\@namedef{PY@tok@w}{\def\PY@tc##1{\textcolor[rgb]{0.73,0.73,0.73}{##1}}}
\@namedef{PY@tok@c}{\let\PY@it=\textit\def\PY@tc##1{\textcolor[rgb]{0.24,0.48,0.48}{##1}}}
\@namedef{PY@tok@cp}{\def\PY@tc##1{\textcolor[rgb]{0.61,0.40,0.00}{##1}}}
\@namedef{PY@tok@k}{\let\PY@bf=\textbf\def\PY@tc##1{\textcolor[rgb]{0.00,0.50,0.00}{##1}}}
\@namedef{PY@tok@kp}{\def\PY@tc##1{\textcolor[rgb]{0.00,0.50,0.00}{##1}}}
\@namedef{PY@tok@kt}{\def\PY@tc##1{\textcolor[rgb]{0.69,0.00,0.25}{##1}}}
\@namedef{PY@tok@o}{\def\PY@tc##1{\textcolor[rgb]{0.40,0.40,0.40}{##1}}}
\@namedef{PY@tok@ow}{\let\PY@bf=\textbf\def\PY@tc##1{\textcolor[rgb]{0.67,0.13,1.00}{##1}}}
\@namedef{PY@tok@nb}{\def\PY@tc##1{\textcolor[rgb]{0.00,0.50,0.00}{##1}}}
\@namedef{PY@tok@nf}{\def\PY@tc##1{\textcolor[rgb]{0.00,0.00,1.00}{##1}}}
\@namedef{PY@tok@nc}{\let\PY@bf=\textbf\def\PY@tc##1{\textcolor[rgb]{0.00,0.00,1.00}{##1}}}
\@namedef{PY@tok@nn}{\let\PY@bf=\textbf\def\PY@tc##1{\textcolor[rgb]{0.00,0.00,1.00}{##1}}}
\@namedef{PY@tok@ne}{\let\PY@bf=\textbf\def\PY@tc##1{\textcolor[rgb]{0.80,0.25,0.22}{##1}}}
\@namedef{PY@tok@nv}{\def\PY@tc##1{\textcolor[rgb]{0.10,0.09,0.49}{##1}}}
\@namedef{PY@tok@no}{\def\PY@tc##1{\textcolor[rgb]{0.53,0.00,0.00}{##1}}}
\@namedef{PY@tok@nl}{\def\PY@tc##1{\textcolor[rgb]{0.46,0.46,0.00}{##1}}}
\@namedef{PY@tok@ni}{\let\PY@bf=\textbf\def\PY@tc##1{\textcolor[rgb]{0.44,0.44,0.44}{##1}}}
\@namedef{PY@tok@na}{\def\PY@tc##1{\textcolor[rgb]{0.41,0.47,0.13}{##1}}}
\@namedef{PY@tok@nt}{\let\PY@bf=\textbf\def\PY@tc##1{\textcolor[rgb]{0.00,0.50,0.00}{##1}}}
\@namedef{PY@tok@nd}{\def\PY@tc##1{\textcolor[rgb]{0.67,0.13,1.00}{##1}}}
\@namedef{PY@tok@s}{\def\PY@tc##1{\textcolor[rgb]{0.73,0.13,0.13}{##1}}}
\@namedef{PY@tok@sd}{\let\PY@it=\textit\def\PY@tc##1{\textcolor[rgb]{0.73,0.13,0.13}{##1}}}
\@namedef{PY@tok@si}{\let\PY@bf=\textbf\def\PY@tc##1{\textcolor[rgb]{0.64,0.35,0.47}{##1}}}
\@namedef{PY@tok@se}{\let\PY@bf=\textbf\def\PY@tc##1{\textcolor[rgb]{0.67,0.36,0.12}{##1}}}
\@namedef{PY@tok@sr}{\def\PY@tc##1{\textcolor[rgb]{0.64,0.35,0.47}{##1}}}
\@namedef{PY@tok@ss}{\def\PY@tc##1{\textcolor[rgb]{0.10,0.09,0.49}{##1}}}
\@namedef{PY@tok@sx}{\def\PY@tc##1{\textcolor[rgb]{0.00,0.50,0.00}{##1}}}
\@namedef{PY@tok@m}{\def\PY@tc##1{\textcolor[rgb]{0.40,0.40,0.40}{##1}}}
\@namedef{PY@tok@gh}{\let\PY@bf=\textbf\def\PY@tc##1{\textcolor[rgb]{0.00,0.00,0.50}{##1}}}
\@namedef{PY@tok@gu}{\let\PY@bf=\textbf\def\PY@tc##1{\textcolor[rgb]{0.50,0.00,0.50}{##1}}}
\@namedef{PY@tok@gd}{\def\PY@tc##1{\textcolor[rgb]{0.63,0.00,0.00}{##1}}}
\@namedef{PY@tok@gi}{\def\PY@tc##1{\textcolor[rgb]{0.00,0.52,0.00}{##1}}}
\@namedef{PY@tok@gr}{\def\PY@tc##1{\textcolor[rgb]{0.89,0.00,0.00}{##1}}}
\@namedef{PY@tok@ge}{\let\PY@it=\textit}
\@namedef{PY@tok@gs}{\let\PY@bf=\textbf}
\@namedef{PY@tok@ges}{\let\PY@bf=\textbf\let\PY@it=\textit}
\@namedef{PY@tok@gp}{\let\PY@bf=\textbf\def\PY@tc##1{\textcolor[rgb]{0.00,0.00,0.50}{##1}}}
\@namedef{PY@tok@go}{\def\PY@tc##1{\textcolor[rgb]{0.44,0.44,0.44}{##1}}}
\@namedef{PY@tok@gt}{\def\PY@tc##1{\textcolor[rgb]{0.00,0.27,0.87}{##1}}}
\@namedef{PY@tok@err}{\def\PY@bc##1{{\setlength{\fboxsep}{\string -\fboxrule}\fcolorbox[rgb]{1.00,0.00,0.00}{1,1,1}{\strut ##1}}}}
\@namedef{PY@tok@kc}{\let\PY@bf=\textbf\def\PY@tc##1{\textcolor[rgb]{0.00,0.50,0.00}{##1}}}
\@namedef{PY@tok@kd}{\let\PY@bf=\textbf\def\PY@tc##1{\textcolor[rgb]{0.00,0.50,0.00}{##1}}}
\@namedef{PY@tok@kn}{\let\PY@bf=\textbf\def\PY@tc##1{\textcolor[rgb]{0.00,0.50,0.00}{##1}}}
\@namedef{PY@tok@kr}{\let\PY@bf=\textbf\def\PY@tc##1{\textcolor[rgb]{0.00,0.50,0.00}{##1}}}
\@namedef{PY@tok@bp}{\def\PY@tc##1{\textcolor[rgb]{0.00,0.50,0.00}{##1}}}
\@namedef{PY@tok@fm}{\def\PY@tc##1{\textcolor[rgb]{0.00,0.00,1.00}{##1}}}
\@namedef{PY@tok@vc}{\def\PY@tc##1{\textcolor[rgb]{0.10,0.09,0.49}{##1}}}
\@namedef{PY@tok@vg}{\def\PY@tc##1{\textcolor[rgb]{0.10,0.09,0.49}{##1}}}
\@namedef{PY@tok@vi}{\def\PY@tc##1{\textcolor[rgb]{0.10,0.09,0.49}{##1}}}
\@namedef{PY@tok@vm}{\def\PY@tc##1{\textcolor[rgb]{0.10,0.09,0.49}{##1}}}
\@namedef{PY@tok@sa}{\def\PY@tc##1{\textcolor[rgb]{0.73,0.13,0.13}{##1}}}
\@namedef{PY@tok@sb}{\def\PY@tc##1{\textcolor[rgb]{0.73,0.13,0.13}{##1}}}
\@namedef{PY@tok@sc}{\def\PY@tc##1{\textcolor[rgb]{0.73,0.13,0.13}{##1}}}
\@namedef{PY@tok@dl}{\def\PY@tc##1{\textcolor[rgb]{0.73,0.13,0.13}{##1}}}
\@namedef{PY@tok@s2}{\def\PY@tc##1{\textcolor[rgb]{0.73,0.13,0.13}{##1}}}
\@namedef{PY@tok@sh}{\def\PY@tc##1{\textcolor[rgb]{0.73,0.13,0.13}{##1}}}
\@namedef{PY@tok@s1}{\def\PY@tc##1{\textcolor[rgb]{0.73,0.13,0.13}{##1}}}
\@namedef{PY@tok@mb}{\def\PY@tc##1{\textcolor[rgb]{0.40,0.40,0.40}{##1}}}
\@namedef{PY@tok@mf}{\def\PY@tc##1{\textcolor[rgb]{0.40,0.40,0.40}{##1}}}
\@namedef{PY@tok@mh}{\def\PY@tc##1{\textcolor[rgb]{0.40,0.40,0.40}{##1}}}
\@namedef{PY@tok@mi}{\def\PY@tc##1{\textcolor[rgb]{0.40,0.40,0.40}{##1}}}
\@namedef{PY@tok@il}{\def\PY@tc##1{\textcolor[rgb]{0.40,0.40,0.40}{##1}}}
\@namedef{PY@tok@mo}{\def\PY@tc##1{\textcolor[rgb]{0.40,0.40,0.40}{##1}}}
\@namedef{PY@tok@ch}{\let\PY@it=\textit\def\PY@tc##1{\textcolor[rgb]{0.24,0.48,0.48}{##1}}}
\@namedef{PY@tok@cm}{\let\PY@it=\textit\def\PY@tc##1{\textcolor[rgb]{0.24,0.48,0.48}{##1}}}
\@namedef{PY@tok@cpf}{\let\PY@it=\textit\def\PY@tc##1{\textcolor[rgb]{0.24,0.48,0.48}{##1}}}
\@namedef{PY@tok@c1}{\let\PY@it=\textit\def\PY@tc##1{\textcolor[rgb]{0.24,0.48,0.48}{##1}}}
\@namedef{PY@tok@cs}{\let\PY@it=\textit\def\PY@tc##1{\textcolor[rgb]{0.24,0.48,0.48}{##1}}}

\def\PYZbs{\char`\\}
\def\PYZus{\char`\_}
\def\PYZob{\char`\{}
\def\PYZcb{\char`\}}
\def\PYZca{\char`\^}
\def\PYZam{\char`\&}
\def\PYZlt{\char`\<}
\def\PYZgt{\char`\>}
\def\PYZsh{\char`\#}
\def\PYZpc{\char`\%}
\def\PYZdl{\char`\$}
\def\PYZhy{\char`\-}
\def\PYZsq{\char`\'}
\def\PYZdq{\char`\"}
\def\PYZti{\char`\~}
% for compatibility with earlier versions
\def\PYZat{@}
\def\PYZlb{[}
\def\PYZrb{]}
\makeatother


    % For linebreaks inside Verbatim environment from package fancyvrb.
    \makeatletter
        \newbox\Wrappedcontinuationbox
        \newbox\Wrappedvisiblespacebox
        \newcommand*\Wrappedvisiblespace {\textcolor{red}{\textvisiblespace}}
        \newcommand*\Wrappedcontinuationsymbol {\textcolor{red}{\llap{\tiny$\m@th\hookrightarrow$}}}
        \newcommand*\Wrappedcontinuationindent {3ex }
        \newcommand*\Wrappedafterbreak {\kern\Wrappedcontinuationindent\copy\Wrappedcontinuationbox}
        % Take advantage of the already applied Pygments mark-up to insert
        % potential linebreaks for TeX processing.
        %        {, <, #, %, $, ' and ": go to next line.
        %        _, }, ^, &, >, - and ~: stay at end of broken line.
        % Use of \textquotesingle for straight quote.
        \newcommand*\Wrappedbreaksatspecials {%
            \def\PYGZus{\discretionary{\char`\_}{\Wrappedafterbreak}{\char`\_}}%
            \def\PYGZob{\discretionary{}{\Wrappedafterbreak\char`\{}{\char`\{}}%
            \def\PYGZcb{\discretionary{\char`\}}{\Wrappedafterbreak}{\char`\}}}%
            \def\PYGZca{\discretionary{\char`\^}{\Wrappedafterbreak}{\char`\^}}%
            \def\PYGZam{\discretionary{\char`\&}{\Wrappedafterbreak}{\char`\&}}%
            \def\PYGZlt{\discretionary{}{\Wrappedafterbreak\char`\<}{\char`\<}}%
            \def\PYGZgt{\discretionary{\char`\>}{\Wrappedafterbreak}{\char`\>}}%
            \def\PYGZsh{\discretionary{}{\Wrappedafterbreak\char`\#}{\char`\#}}%
            \def\PYGZpc{\discretionary{}{\Wrappedafterbreak\char`\%}{\char`\%}}%
            \def\PYGZdl{\discretionary{}{\Wrappedafterbreak\char`\$}{\char`\$}}%
            \def\PYGZhy{\discretionary{\char`\-}{\Wrappedafterbreak}{\char`\-}}%
            \def\PYGZsq{\discretionary{}{\Wrappedafterbreak\textquotesingle}{\textquotesingle}}%
            \def\PYGZdq{\discretionary{}{\Wrappedafterbreak\char`\"}{\char`\"}}%
            \def\PYGZti{\discretionary{\char`\~}{\Wrappedafterbreak}{\char`\~}}%
        }
        % Some characters . , ; ? ! / are not pygmentized.
        % This macro makes them "active" and they will insert potential linebreaks
        \newcommand*\Wrappedbreaksatpunct {%
            \lccode`\~`\.\lowercase{\def~}{\discretionary{\hbox{\char`\.}}{\Wrappedafterbreak}{\hbox{\char`\.}}}%
            \lccode`\~`\,\lowercase{\def~}{\discretionary{\hbox{\char`\,}}{\Wrappedafterbreak}{\hbox{\char`\,}}}%
            \lccode`\~`\;\lowercase{\def~}{\discretionary{\hbox{\char`\;}}{\Wrappedafterbreak}{\hbox{\char`\;}}}%
            \lccode`\~`\:\lowercase{\def~}{\discretionary{\hbox{\char`\:}}{\Wrappedafterbreak}{\hbox{\char`\:}}}%
            \lccode`\~`\?\lowercase{\def~}{\discretionary{\hbox{\char`\?}}{\Wrappedafterbreak}{\hbox{\char`\?}}}%
            \lccode`\~`\!\lowercase{\def~}{\discretionary{\hbox{\char`\!}}{\Wrappedafterbreak}{\hbox{\char`\!}}}%
            \lccode`\~`\/\lowercase{\def~}{\discretionary{\hbox{\char`\/}}{\Wrappedafterbreak}{\hbox{\char`\/}}}%
            \catcode`\.\active
            \catcode`\,\active
            \catcode`\;\active
            \catcode`\:\active
            \catcode`\?\active
            \catcode`\!\active
            \catcode`\/\active
            \lccode`\~`\~
        }
    \makeatother

    \let\OriginalVerbatim=\Verbatim
    \makeatletter
    \renewcommand{\Verbatim}[1][1]{%
        %\parskip\z@skip
        \sbox\Wrappedcontinuationbox {\Wrappedcontinuationsymbol}%
        \sbox\Wrappedvisiblespacebox {\FV@SetupFont\Wrappedvisiblespace}%
        \def\FancyVerbFormatLine ##1{\hsize\linewidth
            \vtop{\raggedright\hyphenpenalty\z@\exhyphenpenalty\z@
                \doublehyphendemerits\z@\finalhyphendemerits\z@
                \strut ##1\strut}%
        }%
        % If the linebreak is at a space, the latter will be displayed as visible
        % space at end of first line, and a continuation symbol starts next line.
        % Stretch/shrink are however usually zero for typewriter font.
        \def\FV@Space {%
            \nobreak\hskip\z@ plus\fontdimen3\font minus\fontdimen4\font
            \discretionary{\copy\Wrappedvisiblespacebox}{\Wrappedafterbreak}
            {\kern\fontdimen2\font}%
        }%

        % Allow breaks at special characters using \PYG... macros.
        \Wrappedbreaksatspecials
        % Breaks at punctuation characters . , ; ? ! and / need catcode=\active
        \OriginalVerbatim[#1,codes*=\Wrappedbreaksatpunct]%
    }
    \makeatother

    % Exact colors from NB
    \definecolor{incolor}{HTML}{303F9F}
    \definecolor{outcolor}{HTML}{D84315}
    \definecolor{cellborder}{HTML}{CFCFCF}
    \definecolor{cellbackground}{HTML}{F7F7F7}

    % prompt
    \makeatletter
    \newcommand{\boxspacing}{\kern\kvtcb@left@rule\kern\kvtcb@boxsep}
    \makeatother
    \newcommand{\prompt}[4]{
        {\ttfamily\llap{{\color{#2}[#3]:\hspace{3pt}#4}}\vspace{-\baselineskip}}
    }
    

    
    % Prevent overflowing lines due to hard-to-break entities
    \sloppy
    % Setup hyperref package
    \hypersetup{
      breaklinks=true,  % so long urls are correctly broken across lines
      colorlinks=true,
      urlcolor=urlcolor,
      linkcolor=linkcolor,
      citecolor=citecolor,
      }
    % Slightly bigger margins than the latex defaults
    
    \geometry{verbose,tmargin=1in,bmargin=1in,lmargin=1in,rmargin=1in}
    
    

\newcounter{none}
\begin{document}
    
    \maketitle
    
    

    
    \begin{tcolorbox}[breakable, size=fbox, boxrule=1pt, pad at break*=1mm,colback=cellbackground, colframe=cellborder]
\prompt{In}{incolor}{ }{\boxspacing}
\begin{Verbatim}[commandchars=\\\{\}]
\PY{k+kn}{import}\PY{+w}{ }\PY{n+nn}{numpy}\PY{+w}{ }\PY{k}{as}\PY{+w}{ }\PY{n+nn}{np}
\PY{k+kn}{import}\PY{+w}{ }\PY{n+nn}{scipy}\PY{n+nn}{.}\PY{n+nn}{stats}\PY{+w}{ }\PY{k}{as}\PY{+w}{ }\PY{n+nn}{st}

\PY{k+kn}{from}\PY{+w}{ }\PY{n+nn}{matplotlib}\PY{+w}{ }\PY{k+kn}{import} \PY{n}{pyplot} \PY{k}{as} \PY{n}{plt}

\PY{k+kn}{import}\PY{+w}{ }\PY{n+nn}{arviz}\PY{+w}{ }\PY{k}{as}\PY{+w}{ }\PY{n+nn}{az}

\PY{k+kn}{from}\PY{+w}{ }\PY{n+nn}{scipy}\PY{+w}{ }\PY{k+kn}{import} \PY{n}{stats}
\PY{k+kn}{from}\PY{+w}{ }\PY{n+nn}{scipy}\PY{n+nn}{.}\PY{n+nn}{special}\PY{+w}{ }\PY{k+kn}{import} \PY{n}{logsumexp}
\PY{k+kn}{from}\PY{+w}{ }\PY{n+nn}{matplotlib}\PY{n+nn}{.}\PY{n+nn}{colors}\PY{+w}{ }\PY{k+kn}{import} \PY{n}{to\PYZus{}hex}

\PY{k+kn}{from}\PY{+w}{ }\PY{n+nn}{src}\PY{+w}{ }\PY{k+kn}{import} \PY{n}{StudentTSource}
\PY{n}{np}\PY{o}{.}\PY{n}{random}\PY{o}{.}\PY{n}{seed}\PY{p}{(}\PY{l+m+mi}{100}\PY{p}{)}
\end{Verbatim}
\end{tcolorbox}

    \section{Ensaio de sensibilidade local --- Separação Cega de Fontes
Bayesiana}\label{ensaio-de-sensibilidade-local-separauxe7uxe3o-cega-de-fontes-bayesiana}

\textbf{Notação geral}

{\def\LTcaptype{none} % do not increment counter
\begin{longtable}[]{@{}
  >{\raggedright\arraybackslash}p{(\linewidth - 2\tabcolsep) * \real{0.5000}}
  >{\raggedright\arraybackslash}p{(\linewidth - 2\tabcolsep) * \real{0.5000}}@{}}
\toprule\noalign{}
\begin{minipage}[b]{\linewidth}\raggedright
Símbolo
\end{minipage} & \begin{minipage}[b]{\linewidth}\raggedright
Significado
\end{minipage} \\
\midrule\noalign{}
\endhead
\bottomrule\noalign{}
\endlastfoot
\(\mathbf{s}(n),\ \mathbf{x}(n)\), \(n=1,\dots,N\) & fontes e misturas;
realização \(\mathcal X = \{\mathbf{x}(n)\}\) \\
\(\underline{\mathbf{A}},\ \underline{\mathbf{B}}\) & matrizes
verdadeiras de mistura e de separação \\
\(\mathbf b = \mathrm{vec}(\mathbf B) \in \mathbb R^{d}\), \(d = L^2\) &
vetorização (\emph{flattening}) de \(\mathbf B\); idem
\(\underline{\mathbf b} = \mathrm{vec}(\underline{\mathbf B})\),
\(\mathbf i = \mathrm{vec}(\mathbf I_2)\) \\
\(p_{s_\ell}(\cdot\,;\xi_s)\) & modelo da \(\ell\)-ésima fonte,
\(\ell = 1,\dots,L\) \\
\(r_s\) & distribuição verdadeira das fontes \\
\(p_{\mathbf b}(\mathbf b;\xi_p)\),
\(p_{\mathbf B\mid\mathcal X}(\mathbf B\mid\mathcal X)\) & distribuições
a priori (sobre \(\mathbf b\)) e a posteriori \\
\(\xi_p = (\boldsymbol\theta_p,\ \sigma_p,\ \nu_p)\) & parâmetros da
priori: moda vetorial (\(\boldsymbol\theta_p\in\mathbb R^d\)), escala e
graus de liberdade \\
\(\xi_s = (\mu_s,\ \sigma_s,\ \nu_s)\) & parâmetros do modelo das fontes
(verossimilhança): localização, escala e graus de liberdade \\
\(\xi_p^{0},\ \xi_s^{0}\) & pontos base; os valores perturbados são
indicados por \(\xi_p^{*}\) e \(\xi_s^{*}\) \\
\(\xi_p^{a}=(\boldsymbol\theta_p^{a},\sigma_p^{a},\nu_p^{a})\),
\(\xi_s^{a}=(\mu_s^{a},\sigma_s^{a},\nu_s^{a})\), \(a\in\{0,*\}\) &
parâmetros base (\(0\)) e perturbados (\(*\)); sem sobrescrito, a
notação é genérica \\
\(\mathbf B_k\), \(\mathbf b_k = \mathrm{vec}(\mathbf B_k)\),
\(k = 1,\dots,K\) & amostras MCMC da posteriori base \\
\end{longtable}
}

\begin{center}\rule{0.5\linewidth}{0.5pt}\end{center}

\subsection{1. Contexto e modelo}\label{contexto-e-modelo}

\textbf{Problema:} estimar \(\mathbf B\) a partir de
\(\mathbf x(n) = \underline{\mathbf A}\,\mathbf s(n)\), com sistema
misturador linear, instantâneo e igualmente determinado (\(L = M\)).

{\def\LTcaptype{none} % do not increment counter
\begin{longtable}[]{@{}
  >{\raggedright\arraybackslash}p{(\linewidth - 2\tabcolsep) * \real{0.5000}}
  >{\raggedright\arraybackslash}p{(\linewidth - 2\tabcolsep) * \real{0.5000}}@{}}
\toprule\noalign{}
\begin{minipage}[b]{\linewidth}\raggedright
Item
\end{minipage} & \begin{minipage}[b]{\linewidth}\raggedright
Valor
\end{minipage} \\
\midrule\noalign{}
\endhead
\bottomrule\noalign{}
\endlastfoot
Fontes / observações & \(L = 2\), \(N = 1000\) \\
Mistura &
\(\underline{\mathbf A} = \tfrac{\sqrt2}{2}\begin{bmatrix}1 & -1\\ 1 & 1\end{bmatrix}\)
(rotação de 45°) \\
Separação ideal &
\(\underline{\mathbf B} = \underline{\mathbf A}^{-1} = \tfrac{\sqrt2}{2}\begin{bmatrix}1 & 1\\ -1 & 1\end{bmatrix}\) \\
Realizações & \(R = 25\) (sementes 0--24, as mesmas para todas as
fontes) \\
\end{longtable}
}

\textbf{Log-posteriori normalizada:}

\[\frac1N\log p_{\mathbf B\mid\mathcal X}(\mathbf B\mid\mathcal X;\xi_p,\xi_s) = \log|\det\mathbf B| + \frac1N\sum_{n=1}^{N}\sum_{\ell=1}^{L}\log p_{s_\ell}\!\big([\mathbf B\mathbf x(n)]_\ell;\xi_s\big) + \frac1N\log p_{\mathbf b}(\mathbf b;\xi_p) + \text{cte}, \qquad \mathbf b = \mathrm{vec}(\mathbf B)\]

\begin{center}\rule{0.5\linewidth}{0.5pt}\end{center}

\subsection{2. Fontes consideradas}\label{fontes-consideradas}

Para as distribuições das fontes, toma-se como forma funcional a
distribuição Student‑t com parâmetros de localização e escala:
\(p_{s_\ell}(s;\xi_s) \sim \mathcal{t}(\mu_s,\ \sigma_s,\ \nu_s)\).
Assim, tem-se que

\[p_{s_\ell}(s;\xi_s) = \frac{\Gamma\!\left(\frac{\nu_s+1}{2}\right)}{\Gamma\!\left(\frac{\nu_s}{2}\right)\sqrt{\nu_s\pi}\,\sigma_s}\left[1 + \frac{1}{\nu_s}\left(\frac{s-\mu_s}{\sigma_s}\right)^{2}\right]^{-\frac{\nu_s+1}{2}}.\]

São consideradas \textbf{dez} distribuições verdadeiras \(r_s\), todas
Student‑t com:

\begin{itemize}
\tightlist
\item
  \(\mu_s^{0} = 0\),
\item
  \((\sigma_s^{0})^2 = (\nu_s^{0}-2)/\nu_s^{0}\),
\item
  \(\nu_s^{0} \in \{3, 4, ..., 12 \}\).
\end{itemize}

Assim, as fontes rs1--rs10 possuem, todas, \textbf{variância unitária},
e diferem apenas na forma (peso de cauda e concentração em torno de
\(\mu_s\)). Uma vez que a Student-t é uma generalização da gaussiana,
(\(\mathcal{t}(\cdot) \rightarrow N(\cdot)\) à medida que
\(\nu \rightarrow \infty\)), a configuração realiza um \textbf{controle
de gaussianidade das fontes}. A aproximação da distribuição normal se
torna progressivamente melhor de rs1 a rs10.

{\def\LTcaptype{none} % do not increment counter
\begin{longtable}[]{@{}
  >{\raggedright\arraybackslash}p{(\linewidth - 20\tabcolsep) * \real{0.0909}}
  >{\raggedright\arraybackslash}p{(\linewidth - 20\tabcolsep) * \real{0.0909}}
  >{\raggedright\arraybackslash}p{(\linewidth - 20\tabcolsep) * \real{0.0909}}
  >{\raggedright\arraybackslash}p{(\linewidth - 20\tabcolsep) * \real{0.0909}}
  >{\raggedright\arraybackslash}p{(\linewidth - 20\tabcolsep) * \real{0.0909}}
  >{\raggedright\arraybackslash}p{(\linewidth - 20\tabcolsep) * \real{0.0909}}
  >{\raggedright\arraybackslash}p{(\linewidth - 20\tabcolsep) * \real{0.0909}}
  >{\raggedright\arraybackslash}p{(\linewidth - 20\tabcolsep) * \real{0.0909}}
  >{\raggedright\arraybackslash}p{(\linewidth - 20\tabcolsep) * \real{0.0909}}
  >{\raggedright\arraybackslash}p{(\linewidth - 20\tabcolsep) * \real{0.0909}}
  >{\raggedright\arraybackslash}p{(\linewidth - 20\tabcolsep) * \real{0.0909}}@{}}
\toprule\noalign{}
\begin{minipage}[b]{\linewidth}\raggedright
\(r_s\)
\end{minipage} & \begin{minipage}[b]{\linewidth}\raggedright
rs1
\end{minipage} & \begin{minipage}[b]{\linewidth}\raggedright
rs2
\end{minipage} & \begin{minipage}[b]{\linewidth}\raggedright
rs3
\end{minipage} & \begin{minipage}[b]{\linewidth}\raggedright
rs4
\end{minipage} & \begin{minipage}[b]{\linewidth}\raggedright
rs5
\end{minipage} & \begin{minipage}[b]{\linewidth}\raggedright
rs6
\end{minipage} & \begin{minipage}[b]{\linewidth}\raggedright
rs7
\end{minipage} & \begin{minipage}[b]{\linewidth}\raggedright
rs8
\end{minipage} & \begin{minipage}[b]{\linewidth}\raggedright
rs9
\end{minipage} & \begin{minipage}[b]{\linewidth}\raggedright
rs10
\end{minipage} \\
\midrule\noalign{}
\endhead
\bottomrule\noalign{}
\endlastfoot
\(\nu_s\) & 3 & 4 & 5 & 6 & 7 & 8 & 9 & 10 & 11 & 12 \\
\(\sigma_s^2\) & 1/3 & 1/2 & 3/5 & 2/3 & 5/7 & 3/4 & 7/9 & 4/5 & 9/11 &
5/6 \\
\end{longtable}
}

Intervalos de Maior Densidade (Highest Density Intervals - HDIs) e
Divergência KL em relação à Gaussiana.

{\def\LTcaptype{none} % do not increment counter
\begin{longtable}[]{@{}
  >{\raggedright\arraybackslash}p{(\linewidth - 6\tabcolsep) * \real{0.2500}}
  >{\raggedright\arraybackslash}p{(\linewidth - 6\tabcolsep) * \real{0.2500}}
  >{\raggedright\arraybackslash}p{(\linewidth - 6\tabcolsep) * \real{0.2500}}
  >{\raggedright\arraybackslash}p{(\linewidth - 6\tabcolsep) * \real{0.2500}}@{}}
\toprule\noalign{}
\begin{minipage}[b]{\linewidth}\raggedright
Distribuição
\end{minipage} & \begin{minipage}[b]{\linewidth}\raggedright
HDI de 95\%
\end{minipage} & \begin{minipage}[b]{\linewidth}\raggedright
Variação do HDI vs.~gaussiana
\end{minipage} & \begin{minipage}[b]{\linewidth}\raggedright
\(D_{KL}(r_s \,\|\, \mathcal{N}(0,1))\)
\end{minipage} \\
\midrule\noalign{}
\endhead
\bottomrule\noalign{}
\endlastfoot
rs1 (\(\nu_s=3\)) & \(\pm\,1{,}837\) & \(-6{,}3\%\) & \(0{,}1948\) \\
rs2 (\(\nu_s=4\)) & \(\pm\,1{,}963\) & \(+0{,}2\%\) & \(0{,}0838\) \\
rs3 (\(\nu_s=5\)) & \(\pm\,1{,}991\) & \(+1{,}6\%\) & \(0{,}0468\) \\
rs4 (\(\nu_s=6\)) & \(\pm\,1{,}998\) & \(+1{,}9\%\) & \(0{,}0299\) \\
rs5 (\(\nu_s=7\)) & \(\pm\,1{,}998\) & \(+2{,}0\%\) & \(0{,}0208\) \\
rs6 (\(\nu_s=8\)) & \(\pm\,1{,}997\) & \(+1{,}9\%\) & \(0{,}0153\) \\
rs7 (\(\nu_s=9\)) & \(\pm\,1{,}995\) & \(+1{,}8\%\) & \(0{,}0117\) \\
rs8 (\(\nu_s=10\)) & \(\pm\,1{,}993\) & \(+1{,}7\%\) & \(0{,}0092\) \\
rs9 (\(\nu_s=11\)) & \(\pm\,1{,}991\) & \(+1{,}6\%\) & \(0{,}0075\) \\
rs10 (\(\nu_s=12\)) & \(\pm\,1{,}989\) & \(+1{,}5\%\) & \(0{,}0062\) \\
Gaussiana \(\mathcal{N}(0,1)\) & \(\pm\,1{,}960\) & \(0{,}0\%\) &
\(0\) \\
\end{longtable}
}

    \begin{tcolorbox}[breakable, size=fbox, boxrule=1pt, pad at break*=1mm,colback=cellbackground, colframe=cellborder]
\prompt{In}{incolor}{5}{\boxspacing}
\begin{Verbatim}[commandchars=\\\{\}]
\PY{c+c1}{\PYZsh{}\PYZsh{}\PYZsh{}\PYZsh{}\PYZsh{}\PYZsh{}\PYZsh{}\PYZsh{}\PYZsh{}\PYZsh{}\PYZsh{}\PYZsh{}\PYZsh{}\PYZsh{}\PYZsh{}}
\PY{c+c1}{\PYZsh{} Plot fontes \PYZsh{}}
\PY{c+c1}{\PYZsh{}\PYZsh{}\PYZsh{}\PYZsh{}\PYZsh{}\PYZsh{}\PYZsh{}\PYZsh{}\PYZsh{}\PYZsh{}\PYZsh{}\PYZsh{}\PYZsh{}\PYZsh{}\PYZsh{}}

\PY{c+c1}{\PYZsh{} Graus de liberdade e valores x avaliados}
\PY{n}{evaluated\PYZus{}dfs} \PY{o}{=} \PY{n+nb}{range}\PY{p}{(}\PY{l+m+mi}{3}\PY{p}{,}\PY{l+m+mi}{13}\PY{p}{)}
\PY{n}{evaluated\PYZus{}x} \PY{o}{=} \PY{n}{np}\PY{o}{.}\PY{n}{linspace}\PY{p}{(}\PY{o}{\PYZhy{}}\PY{l+m+mi}{5}\PY{p}{,}\PY{l+m+mi}{5}\PY{p}{,}\PY{l+m+mi}{1000}\PY{p}{)}

\PY{c+c1}{\PYZsh{} Colormap (mesmo padrão dos métodos de IS)}
\PY{n}{cmap} \PY{o}{=} \PY{n}{plt}\PY{o}{.}\PY{n}{get\PYZus{}cmap}\PY{p}{(}\PY{l+s+s1}{\PYZsq{}}\PY{l+s+s1}{magma}\PY{l+s+s1}{\PYZsq{}}\PY{p}{)}
\PY{n}{colors} \PY{o}{=} \PY{n}{cmap}\PY{p}{(}\PY{n}{np}\PY{o}{.}\PY{n}{linspace}\PY{p}{(}\PY{l+m+mf}{0.0}\PY{p}{,} \PY{l+m+mf}{0.8}\PY{p}{,} \PY{n+nb}{len}\PY{p}{(}\PY{n}{evaluated\PYZus{}dfs}\PY{p}{)}\PY{p}{)}\PY{p}{)}

\PY{c+c1}{\PYZsh{} Instantiate sources}
\PY{n}{sources} \PY{o}{=} \PY{p}{\PYZob{}}\PY{n+nb}{str}\PY{p}{(}\PY{n}{df}\PY{p}{)}\PY{p}{:} \PY{n}{StudentTSource}\PY{p}{(}
    \PY{n}{mu}\PY{o}{=}\PY{l+m+mi}{0}\PY{p}{,}
    \PY{n}{sigma}\PY{o}{=}\PY{n}{np}\PY{o}{.}\PY{n}{sqrt}\PY{p}{(}\PY{p}{(}\PY{n}{df}\PY{o}{\PYZhy{}}\PY{l+m+mi}{2}\PY{p}{)}\PY{o}{/}\PY{n}{df}\PY{p}{)}\PY{p}{,}
    \PY{n}{df}\PY{o}{=}\PY{n}{df}
\PY{p}{)} \PY{k}{for} \PY{n}{df} \PY{o+ow}{in} \PY{n}{evaluated\PYZus{}dfs}\PY{p}{\PYZcb{}}

\PY{n}{fig}\PY{p}{,} \PY{n}{ax} \PY{o}{=} \PY{n}{plt}\PY{o}{.}\PY{n}{subplots}\PY{p}{(}
    \PY{n}{nrows}\PY{o}{=}\PY{l+m+mi}{1}\PY{p}{,} \PY{n}{ncols}\PY{o}{=}\PY{l+m+mi}{1}\PY{p}{,}
    \PY{n}{figsize}\PY{o}{=}\PY{p}{(}\PY{l+m+mi}{20}\PY{p}{,}\PY{l+m+mi}{7}\PY{p}{)}
\PY{p}{)}

\PY{n}{ax}\PY{o}{.}\PY{n}{set\PYZus{}xlabel}\PY{p}{(}\PY{l+s+s1}{\PYZsq{}}\PY{l+s+s1}{x}\PY{l+s+s1}{\PYZsq{}}\PY{p}{,} \PY{n}{fontsize}\PY{o}{=}\PY{l+m+mi}{15}\PY{p}{)}
\PY{n}{ax}\PY{o}{.}\PY{n}{set\PYZus{}ylabel}\PY{p}{(}\PY{l+s+s1}{\PYZsq{}}\PY{l+s+s1}{density}\PY{l+s+s1}{\PYZsq{}}\PY{p}{,} \PY{n}{fontsize}\PY{o}{=}\PY{l+m+mi}{15}\PY{p}{)}

\PY{c+c1}{\PYZsh{} Plot t}
\PY{k}{for} \PY{n}{color}\PY{p}{,} \PY{p}{(}\PY{n}{df}\PY{p}{,} \PY{n}{s}\PY{p}{)} \PY{o+ow}{in} \PY{n+nb}{zip}\PY{p}{(}\PY{n}{colors}\PY{p}{,} \PY{n}{sources}\PY{o}{.}\PY{n}{items}\PY{p}{(}\PY{p}{)}\PY{p}{)}\PY{p}{:}
    \PY{n}{ax}\PY{o}{.}\PY{n}{plot}\PY{p}{(}
        \PY{n}{evaluated\PYZus{}x}\PY{p}{,}
        \PY{n}{s}\PY{o}{.}\PY{n}{get}\PY{p}{(}\PY{p}{)}\PY{p}{(}\PY{n}{evaluated\PYZus{}x}\PY{p}{)}\PY{p}{,}
        \PY{n}{color}\PY{o}{=}\PY{n}{color}\PY{p}{,}
        \PY{n}{label}\PY{o}{=}\PY{l+s+s1}{\PYZsq{}}\PY{l+s+s1}{rs}\PY{l+s+si}{\PYZob{}\PYZcb{}}\PY{l+s+s1}{ (df=}\PY{l+s+si}{\PYZob{}\PYZcb{}}\PY{l+s+s1}{)}\PY{l+s+s1}{\PYZsq{}}\PY{o}{.}\PY{n}{format}\PY{p}{(}\PY{n+nb}{int}\PY{p}{(}\PY{n}{df}\PY{p}{)}\PY{o}{\PYZhy{}}\PY{l+m+mi}{2}\PY{p}{,} \PY{n}{df}\PY{p}{)}
    \PY{p}{)}

\PY{c+c1}{\PYZsh{} Plot gaussiana N(0,1) com destaque}
\PY{n}{ax}\PY{o}{.}\PY{n}{plot}\PY{p}{(}
    \PY{n}{evaluated\PYZus{}x}\PY{p}{,}
    \PY{n}{st}\PY{o}{.}\PY{n}{norm}\PY{o}{.}\PY{n}{pdf}\PY{p}{(}\PY{n}{evaluated\PYZus{}x}\PY{p}{,} \PY{n}{loc}\PY{o}{=}\PY{l+m+mi}{0}\PY{p}{,} \PY{n}{scale}\PY{o}{=}\PY{l+m+mi}{1}\PY{p}{)}\PY{p}{,}
    \PY{n}{color}\PY{o}{=}\PY{n}{cmap}\PY{p}{(}\PY{l+m+mf}{0.88}\PY{p}{)}\PY{p}{,}
    \PY{n}{linestyle}\PY{o}{=}\PY{l+s+s1}{\PYZsq{}}\PY{l+s+s1}{\PYZhy{}\PYZhy{}}\PY{l+s+s1}{\PYZsq{}}\PY{p}{,}
    \PY{n}{linewidth}\PY{o}{=}\PY{l+m+mi}{3}\PY{p}{,}
    \PY{n}{label}\PY{o}{=}\PY{l+s+sa}{r}\PY{l+s+s1}{\PYZsq{}}\PY{l+s+s1}{Gaussiana \PYZdl{}}\PY{l+s+s1}{\PYZbs{}}\PY{l+s+s1}{mathcal}\PY{l+s+si}{\PYZob{}N\PYZcb{}}\PY{l+s+s1}{(0,1)\PYZdl{}}\PY{l+s+s1}{\PYZsq{}}\PY{p}{,}
    \PY{n}{zorder}\PY{o}{=}\PY{l+m+mi}{10}
\PY{p}{)}

\PY{n}{l} \PY{o}{=} \PY{n}{ax}\PY{o}{.}\PY{n}{legend}\PY{p}{(}\PY{n}{fontsize}\PY{o}{=}\PY{l+m+mi}{15}\PY{p}{)}
\end{Verbatim}
\end{tcolorbox}

    \begin{center}
    \adjustimage{max size={0.9\linewidth}{0.9\paperheight}}{notebook_corrigido_files/notebook_corrigido_2_0.png}
    \end{center}
    { \hspace*{\fill} \\}
    
    \subsection{3. Priori considerada}\label{priori-considerada}

Para a distribuição a priori, adota-se como forma funcional distribuição
Student-t multivariada sobre
\(\mathbf b = \mathrm{vec}(\mathbf B) \in \mathbb R^{d}\), com
\(d = L^2\) e matriz de covariância \(\sigma_p^2\mathbf I_d\). Assim,
tem-se
\(p_{\mathbf b}(\mathbf b;\xi_p) \sim \mathcal{T}(\boldsymbol\theta_p,\sigma_p,\nu_p)\):

\[p_{\mathbf b}(\mathbf b;\xi_p) = \frac{\Gamma\!\left(\frac{\nu_p+d}{2}\right)}{\Gamma\!\left(\frac{\nu_p}{2}\right)(\nu_p\pi)^{d/2}\,\sigma_p^{d}}\left[1 + \frac{(\mathbf b-\boldsymbol\theta_p)^\top(\mathbf b-\boldsymbol\theta_p)}{\nu_p\,\sigma_p^{2}}\right]^{-\frac{\nu_p+d}{2}}\]

É a generalização, com cauda pesada, da priori gaussiana usada na
Qualificação.

\begin{center}\rule{0.5\linewidth}{0.5pt}\end{center}

\subsection{4. Objetivos científicos e
perguntas}\label{objetivos-cientuxedficos-e-perguntas}

\textbf{Objetivo:} quantificar o quanto \(p_{\mathbf B\mid\mathcal X}\)
é sensível a perturbações locais de \(\xi_p\) e \(\xi_s\).

\textbf{Perguntas:}

\begin{enumerate}
\def\labelenumi{\arabic{enumi}.}
\tightlist
\item
  \textbf{Priori:} quanto os momentos a posteriori de \(b_{ij}\) (média,
  variância, assimetria, curtose) mudam com perturbações em
  \(\boldsymbol\theta_p\), \(\sigma_p\) e \(\nu_p\)?
\item
  \textbf{Modelo das fontes:} quanto os momentos a posteriori de
  \(b_{ij}\) (média, variância, assimetria, curtose) mudam com
  perturbações em \(\mu_s\), \(\sigma_s\) e \(\nu_s\)?
\item
  \textbf{Efeito de \(r_s\):} a sensibilidade muda ao passar de fontes
  de cauda pesada (rs1) para fontes quase gaussianas (rs10)?
\item
  \textbf{Priori × verossimilhança:} com \(N = 1000\), qual dos dois
  componentes domina a posteriori?
\end{enumerate}

\begin{center}\rule{0.5\linewidth}{0.5pt}\end{center}

\subsection{5. Importance sampling
autonormalizado}\label{importance-sampling-autonormalizado}

A amostragem de importância permite estimar esperanças e momentos sob
uma distribuição alvo a partir de amostras de outra distribuição, mais
simples de amostrar ou já disponível. A mesma amostra pode ser
reutilizada para diferentes funções \(h\) e diferentes alvos, ajustando
os \textbf{pesos de importância}.

\textbf{Ideia geral e derivação (Paulino et al., 2018, p.~290--291, eq.
7.18 e 7.19):} sejam \(p^{0}(\boldsymbol{\theta})\) a densidade base e
\(p^{\ast}(\boldsymbol{\theta})\) a densidade alvo, definidas sobre o
espaço paramétrico \(\Phi_{\boldsymbol{\theta}}\), com \(p^{0}>0\). Para
uma função integrável \(h\):

\[\begin{aligned}
\mathcal{H}^{\ast}
&= \mathbb{E}_{p^{\ast}}[h(\boldsymbol{\theta})] \\
&= \int_{\Phi_{\boldsymbol{\theta}}} h(\boldsymbol{\theta})\,p^{\ast}(\boldsymbol{\theta})\,\mathrm{d}\boldsymbol{\theta} \\
&= \int_{\Phi_{\boldsymbol{\theta}}} h(\boldsymbol{\theta})\,\underbrace{\frac{p^{\ast}(\boldsymbol{\theta})}{p^{0}(\boldsymbol{\theta})}}_{w(\boldsymbol{\theta})}\,p^{0}(\boldsymbol{\theta})\,\mathrm{d}\boldsymbol{\theta} \\
&= \mathbb{E}_{p^{0}}[h(\boldsymbol{\theta})\,w(\boldsymbol{\theta})].
\end{aligned}\]

Com \(\boldsymbol{\theta}_k \sim p^{0}\) e
\(w_k = w(\boldsymbol{\theta}_k)\), a estimativa de Monte Carlo é:

\[\widehat{\mathcal{H}}^{\ast}_{\mathrm{IS}} = \frac{1}{K}\sum_{k=1}^{K} h(\boldsymbol{\theta}_k)\,w_k.\]

Na forma autonormalizada, os pesos são divididos pela sua soma:

\[\widehat{\mathcal{H}}^{\ast}_{\mathrm{IS}}
= \frac{\sum_{k=1}^{K} h(\boldsymbol{\theta}_k)\,w_k}{\sum_{k=1}^{K} w_k}
= \sum_{k=1}^{K} \tilde{w}_k\, h(\boldsymbol{\theta}_k),
\qquad
\tilde{w}_k = \frac{w_k}{\sum_{k'=1}^{K} w_{k'}}.\]

\textbf{Aplicação ao ensaio:} amostras
\(\mathbf{B}_k \sim p_{\mathbf{B}\mid\mathcal{X}}(\mathbf{B}\mid\mathcal{X};\xi_p^{0},\xi_s^{0})\),
com \(k=1,\dots,K\):

\[\widehat{\mathbb{E}}_{\xi_p^{\ast},\xi_s^{\ast}}\big[h(\mathbf{B})\big]
= \sum_{k=1}^{K} \tilde{w}_k\, h(\mathbf{B}_k),
\qquad
\tilde{w}_k = \frac{w_k}{\sum_{k'=1}^{K} w_{k'}}.\]

\textbf{Pesos brutos \(w_k\):}

\emph{Posteriori:}

\[\log w_k
= \log p_{\mathbf{B}\mid\mathcal{X}}(\mathbf{B}_k\mid\mathcal{X};\xi_p^{\ast},\xi_s^{\ast})
- \log p_{\mathbf{B}\mid\mathcal{X}}(\mathbf{B}_k\mid\mathcal{X};\xi_p^{0},\xi_s^{0}).\]

\emph{Priori}, quando \(\xi_s^{\ast}=\xi_s^{0}\) (pesos equivalentes
após a autonormalização):

\[\log w_k = \log p_{\mathbf{b}}(\mathbf{b}_k;\xi_p^{\ast}) - \log p_{\mathbf{b}}(\mathbf{b}_k;\xi_p^{0}).\]

\emph{Modelo das fontes}, quando \(\xi_p^{\ast}=\xi_p^{0}\) (pesos
equivalentes após a autonormalização):

\[\log w_k
= \sum_{n=1}^{N}\sum_{\ell=1}^{L}
\Big[
\log p_{s_\ell}\big([\mathbf{B}_k\mathbf{x}(n)]_\ell;\xi_s^{\ast}\big)
- \log p_{s_\ell}\big([\mathbf{B}_k\mathbf{x}(n)]_\ell;\xi_s^{0}\big)
\Big].\]

O termo \(N\log|\det\mathbf{B}_k|\) se cancela.

A soma das diferenças de log-priori e log-verossimilhança fornece os
pesos para alterações simultâneas de \(\xi_p\) e \(\xi_s\), usados na
estimativa autonormalizada.

\textbf{Quantidades de interesse:} média, variância, assimetria e
curtose de cada \(b_{ij}\), ponderadas por \(\tilde{w}_k\). A média
ponderada é a estimativa \(\widehat{\mathbf{B}}_{\mathrm{MMSE}}^{\ast}\)
sob \((\xi_p^{\ast},\xi_s^{\ast})\).

\begin{center}\rule{0.5\linewidth}{0.5pt}\end{center}

\subsection{6. Parâmetros perturbados e
intervalos}\label{paruxe2metros-perturbados-e-intervalos}

\subsubsection{\texorpdfstring{6.1 Variação da priori
(\(\xi_s^{*} = \xi_s^{0}=r_{s}\), para cada uma das fontes
rs1-rs10)}{6.1 Variação da priori (\textbackslash xi\_s\^{}\{*\} = \textbackslash xi\_s\^{}\{0\}=r\_\{s\}, para cada uma das fontes rs1-rs10)}}\label{variauxe7uxe3o-da-priori-xi_s-xi_s0r_s-para-cada-uma-das-fontes-rs1-rs10}

{\def\LTcaptype{none} % do not increment counter
\begin{longtable}[]{@{}
  >{\raggedright\arraybackslash}p{(\linewidth - 6\tabcolsep) * \real{0.2500}}
  >{\raggedright\arraybackslash}p{(\linewidth - 6\tabcolsep) * \real{0.2500}}
  >{\raggedright\arraybackslash}p{(\linewidth - 6\tabcolsep) * \real{0.2500}}
  >{\raggedright\arraybackslash}p{(\linewidth - 6\tabcolsep) * \real{0.2500}}@{}}
\toprule\noalign{}
\begin{minipage}[b]{\linewidth}\raggedright
Perturbação
\end{minipage} & \begin{minipage}[b]{\linewidth}\raggedright
Base
\end{minipage} & \begin{minipage}[b]{\linewidth}\raggedright
Alvos
\end{minipage} & \begin{minipage}[b]{\linewidth}\raggedright
Justificativa do intervalo
\end{minipage} \\
\midrule\noalign{}
\endhead
\bottomrule\noalign{}
\endlastfoot
\textbf{Moda} \(\boldsymbol\theta_p\) &
\(\boldsymbol\theta_p^{0}=\tfrac12(\underline{\mathbf b}+\mathbf i)\),
\(\sigma_p^{0}=\sqrt{1/3}\), \(\nu_p^{0}=3\) &
\(\boldsymbol\theta_p^{*}(\alpha)=(1-\alpha)\underline{\mathbf b} + \alpha\,\mathbf i\),
\(\alpha\in[0,1]\) (50 pts) & Homotopia da moda da priori, indo da
solução ideal à identidade; a base no meio do caminho mantém as duas
pontas dentro da faixa de confiabilidade do PSIS. \\
\textbf{Escala} \(\sigma_p\) &
\(\boldsymbol\theta_p^{0}=\underline{\mathbf b}\),
\(\sigma_p^{0} = 0{,}5\), \(\nu_p^{0}=3\) &
\(\sigma_p^{*} \in [10^{-1}, 10^{1}]\) (50 pts, log) & Vai de priori
informativa a pouco informativa, utilizando os mesmos extremos 0,1 e 10
do experimento da Qualificação; mantém a variação paramétrica dentro da
faixa permitida pelo \(\hat k\) do PSIS. \\
\textbf{Cauda} \(\nu_p\) &
\(\boldsymbol\theta_p^{0}=\underline{\mathbf b}\),
\(\sigma_p^{0}=\sqrt{1/3}\), \(\nu_p^{0}=3\) &
\(\nu_p^{*} \in \{3,\dots,29\}\) & Vai de cauda pesada até perto da
normal, faixa em que os pesos seguem confiáveis segundo o diagnóstico do
PSIS. \\
\end{longtable}
}

\subsubsection{\texorpdfstring{6.2 Variação do modelo das fontes (priori
fixa: \(\boldsymbol\theta_p^{0}=\underline{\mathbf b}\),
\(\sigma_p^{0}=\sqrt{1/3}\),
\(\nu_p^{0}=3\))}{6.2 Variação do modelo das fontes (priori fixa: \textbackslash boldsymbol\textbackslash theta\_p\^{}\{0\}=\textbackslash underline\{\textbackslash mathbf b\}, \textbackslash sigma\_p\^{}\{0\}=\textbackslash sqrt\{1/3\}, \textbackslash nu\_p\^{}\{0\}=3)}}\label{variauxe7uxe3o-do-modelo-das-fontes-priori-fixa-boldsymboltheta_p0underlinemathbf-b-sigma_p0sqrt13-nu_p03}

{\def\LTcaptype{none} % do not increment counter
\begin{longtable}[]{@{}
  >{\raggedright\arraybackslash}p{(\linewidth - 6\tabcolsep) * \real{0.2500}}
  >{\raggedright\arraybackslash}p{(\linewidth - 6\tabcolsep) * \real{0.2500}}
  >{\raggedright\arraybackslash}p{(\linewidth - 6\tabcolsep) * \real{0.2500}}
  >{\raggedright\arraybackslash}p{(\linewidth - 6\tabcolsep) * \real{0.2500}}@{}}
\toprule\noalign{}
\begin{minipage}[b]{\linewidth}\raggedright
Perturbação
\end{minipage} & \begin{minipage}[b]{\linewidth}\raggedright
Base
\end{minipage} & \begin{minipage}[b]{\linewidth}\raggedright
Alvos
\end{minipage} & \begin{minipage}[b]{\linewidth}\raggedright
Justificativa do intervalo
\end{minipage} \\
\midrule\noalign{}
\endhead
\bottomrule\noalign{}
\endlastfoot
\textbf{Localização} \(\mu_s\) & \(\mu_s^{0} = 0\) &
\(\mu_s^{*} \in [-0{,}2;\ 0{,}2]\) (50 pts) & Faixa ao redor da moda do
modelo das fontes base que é indicada como confiável pelo \(\hat k\) do
PSIS. \\
\textbf{Escala} \(\sigma_s\) & \(\sigma_s^{0}\) de \(r_s\) &
\(\sigma_s^{*}=\sigma_s^{0}\,e^{u}\), \(u\in[-0{,}015;\ 0{,}015]\)
(±1,5\%, 50 pts) & Faixa relativa única, a maior que mantém o \(\hat k\)
do PSIS aceitável para todas as fontes (rs1--rs10). \\
\textbf{Cauda} \(\nu_s\) & \(\nu_s^{0}\) de \(r_s\) &
\(\nu_s^{*} \in \{3,\dots,29\}\), \(\sigma_s^{*}=\sigma_s^{0}\) fixo &
Mesma faixa usada na priori, para comparação direta; a confiabilidade é
checada pelo \(\hat k\) em cada ponto. \\
\end{longtable}
}

\begin{quote}
Na cauda, \(\sigma_s^{*}=\sigma_s^{0}\) fica fixo pela mesma lógica
adotada para \(\sigma_p^{*}=\sigma_p^{0}\) na priori.
\end{quote}

\begin{center}\rule{0.5\linewidth}{0.5pt}\end{center}

\subsection{7. Metodologia --- fluxo
geral}\label{metodologia-fluxo-geral}

Fluxo inspirado no \emph{power-scaling} (Kallioinen et al.):
\textbf{perturbar priori e verossimilhança separadamente avaliar
impactos sobre posteriori-base}. A diferença é que, em vez de usar
potências da priori ou da verossimilhança, perturbam-se individualmente
os parâmetros em (\(\xi_p\), \(\xi_s\)).

\begin{verbatim}
MCMC na posteriori base (ξp⁰ , ξs⁰)
        │
        ├─► Varredura da vizinhança de parâmetros: pesos de IS → PSIS → k̂ → momentos ponderados
        │
        └─► Sensibilidade local: Cov(h(B), η) no ponto base (sem reponderar)
                │
        Comparação: inclinação local × curva de variação local (por fonte, por realização)
\end{verbatim}

\subsubsection{7.1 Amostragem base --- Barker MH
(NumPyro)}\label{amostragem-base-barker-mh-numpyro}

Substitui o Metropolis‑Hastings de passeio aleatório da Qualificação: a
proposta usa o gradiente da log-posteriori.

{\def\LTcaptype{none} % do not increment counter
\begin{longtable}[]{@{}
  >{\raggedright\arraybackslash}p{(\linewidth - 2\tabcolsep) * \real{0.5000}}
  >{\raggedright\arraybackslash}p{(\linewidth - 2\tabcolsep) * \real{0.5000}}@{}}
\toprule\noalign{}
\begin{minipage}[b]{\linewidth}\raggedright
Configuração
\end{minipage} & \begin{minipage}[b]{\linewidth}\raggedright
Valor
\end{minipage} \\
\midrule\noalign{}
\endhead
\bottomrule\noalign{}
\endlastfoot
Cadeias / amostras & 4 cadeias, \(K = 20.000\) amostras finais (5.000
por cadeia) \\
Aquecimento & 5.000 iterações, aceitação‑alvo de 0,5 \\
Convergência & \(\hat R < 1{,}01\) em todos os \(b_{ij}\), avaliado a
cada 500 iterações, 2 avaliações seguidas \\
Máx. iterações & 15.000 \\
Inicialização & \(\underline{\mathbf B} + \boldsymbol\varepsilon\),
\(\varepsilon_{ij} \sim t_3(0;\,0{,}05)\) \\
\end{longtable}
}

\begin{center}\rule{0.5\linewidth}{0.5pt}\end{center}

\subsection{8. Equacionamento}\label{equacionamento}

\subsubsection{\texorpdfstring{8.1 PSIS e diagnóstico \(\hat k\)
(Vehtari et al.,
2024)}{8.1 PSIS e diagnóstico \textbackslash hat k (Vehtari et al., 2024)}}\label{psis-e-diagnuxf3stico-hat-k-vehtari-et-al.-2024}

\textbf{Motivação:} quando a posteriori alvo e a base diferem muito em
certas regiões, a distribuição dos pesos \(w_k\) ganha uma cauda longa.
Poucas amostras passam a dominar a estimativa, que fica instável. O PSIS
controla essa cauda.

\begin{itemize}
\tightlist
\item
  Ajusta uma Pareto generalizada às
  \(M_{\text{cauda}} = \lfloor\min(0{,}2K,\ 3\sqrt K)\rfloor\) maiores
  razões.
\item
  Substitui essas razões pelas estatísticas de ordem esperadas:
  \(w_{(K-M_{\text{cauda}}+m)} = F^{-1}\!\left(\frac{m-1/2}{M_{\text{cauda}}}\right)\),
  \(m=1,\dots,M_{\text{cauda}}\). As estatísticas de ordem são os
  valores esperados, sobre os pesos ordenados, da distribuição Pareto
  ajustada.
\item
  \(\hat k\) = forma da Pareto ajustada, que indica quantos momentos dos
  pesos são finitos. \textbf{\(\hat k < 0{,}7\) → confiável} (com
  \(K = 20.000\), o limiar é
  \(\min(1 - 1/\log_{10}K,\ 0{,}7) = 0{,}7\)).
\item
  Eficiência relativa do MCMC: \(R_{\text{eff}} = \text{ESS}_{\min}/K\).
\end{itemize}

→ A curva \(\hat k\) × parâmetro delimita \textbf{a faixa em que a
varredura da vizinhança dos parâmetros perturbados vale}.

Quando os pesos de importância têm cauda pesada, poucos valores extremos
podem dominar a estimativa, tornando-a muito variável entre realizações.
O PSIS substitui os maiores pesos por quantis da Pareto ajustada,
reduzindo a variabilidade amostral (ruído amostral) desses extremos e,
consequentemente, da estimativa. O objetivo não é diminuir todos os
pesos grandes, mas regularizar sua contribuição, com baixo viés, para
obter estimativas mais estáveis.

\begin{center}\rule{0.5\linewidth}{0.5pt}\end{center}

\textbf{Justificativa da expressão das estatísticas de ordem.}

Seja \(F\) a CDF da Pareto ajustada à cauda. Se \(Y\) segue essa
distribuição, a transformação \(U=F(Y)\) segue uma Uniforme\((0,1)\):

\[P(U\le p)
=
P\big(Y\le F^{-1}(p)\big)
=
p,
\qquad 0<p<1.\]

Como \(F\) é crescente, preserva a ordenação dos \(M_{\text{cauda}}\)
valores:

\[Y_{(m)}=F^{-1}\big(U_{(m)}\big).\]

O PSIS usa uma aproximação baseada em percentis igualmente espaçados,
nos pontos médios de \(M_{\text{cauda}}\) faixas de probabilidade:

\[p_m=\frac{m-1/2}{M_{\text{cauda}}} = F^{-1}\!\left(\frac{m-1/2}{M_{\text{cauda}}}\right), m=1,\dots,M_{\text{cauda}}.\]

Aplicando a inversa da Pareto, esses percentis tornam-se os pesos
suavizados:

\[w_{(K-M_{\text{cauda}}+m)}
=
F^{-1}(p_m).\]

Assim, cada peso observado é substituído pelo quantil da Pareto
associado à sua posição na cauda.

    \begin{tcolorbox}[breakable, size=fbox, boxrule=1pt, pad at break*=1mm,colback=cellbackground, colframe=cellborder]
\prompt{In}{incolor}{26}{\boxspacing}
\begin{Verbatim}[commandchars=\\\{\}]
\PY{c+c1}{\PYZsh{} \PYZhy{}\PYZhy{}\PYZhy{}\PYZhy{}\PYZhy{}\PYZhy{}\PYZhy{}\PYZhy{}\PYZhy{}\PYZhy{}\PYZhy{}\PYZhy{}\PYZhy{}\PYZhy{}\PYZhy{}\PYZhy{}\PYZhy{}\PYZhy{}\PYZhy{}\PYZhy{}\PYZhy{}\PYZhy{}\PYZhy{}\PYZhy{}\PYZhy{}\PYZhy{}\PYZhy{}\PYZhy{}\PYZhy{}\PYZhy{}\PYZhy{}\PYZhy{}\PYZhy{}\PYZhy{}\PYZhy{}\PYZhy{}\PYZhy{}\PYZhy{}\PYZhy{}\PYZhy{}\PYZhy{}\PYZhy{}\PYZhy{}\PYZhy{}\PYZhy{}\PYZhy{}\PYZhy{}\PYZhy{}\PYZhy{}\PYZhy{}\PYZhy{}\PYZhy{}\PYZhy{}\PYZhy{}\PYZhy{}\PYZhy{}\PYZhy{}\PYZhy{}\PYZhy{}\PYZhy{}\PYZhy{}\PYZhy{}\PYZhy{}\PYZhy{}\PYZhy{}\PYZhy{}\PYZhy{}\PYZhy{}\PYZhy{}\PYZhy{}\PYZhy{}\PYZhy{}\PYZhy{}}
\PY{c+c1}{\PYZsh{} Configurações}
\PY{c+c1}{\PYZsh{} \PYZhy{}\PYZhy{}\PYZhy{}\PYZhy{}\PYZhy{}\PYZhy{}\PYZhy{}\PYZhy{}\PYZhy{}\PYZhy{}\PYZhy{}\PYZhy{}\PYZhy{}\PYZhy{}\PYZhy{}\PYZhy{}\PYZhy{}\PYZhy{}\PYZhy{}\PYZhy{}\PYZhy{}\PYZhy{}\PYZhy{}\PYZhy{}\PYZhy{}\PYZhy{}\PYZhy{}\PYZhy{}\PYZhy{}\PYZhy{}\PYZhy{}\PYZhy{}\PYZhy{}\PYZhy{}\PYZhy{}\PYZhy{}\PYZhy{}\PYZhy{}\PYZhy{}\PYZhy{}\PYZhy{}\PYZhy{}\PYZhy{}\PYZhy{}\PYZhy{}\PYZhy{}\PYZhy{}\PYZhy{}\PYZhy{}\PYZhy{}\PYZhy{}\PYZhy{}\PYZhy{}\PYZhy{}\PYZhy{}\PYZhy{}\PYZhy{}\PYZhy{}\PYZhy{}\PYZhy{}\PYZhy{}\PYZhy{}\PYZhy{}\PYZhy{}\PYZhy{}\PYZhy{}\PYZhy{}\PYZhy{}\PYZhy{}\PYZhy{}\PYZhy{}\PYZhy{}\PYZhy{}}

\PY{n}{n\PYZus{}curves} \PY{o}{=} \PY{l+m+mi}{10}
\PY{n}{x} \PY{o}{=} \PY{n}{np}\PY{o}{.}\PY{n}{linspace}\PY{p}{(}\PY{o}{\PYZhy{}}\PY{l+m+mf}{0.5}\PY{p}{,} \PY{l+m+mi}{10}\PY{p}{,} \PY{l+m+mi}{2000}\PY{p}{)}

\PY{c+c1}{\PYZsh{} Valores de referência (mantidos fixos quando outro parâmetro varia)}
\PY{n}{u\PYZus{}ref}\PY{p}{,} \PY{n}{k\PYZus{}ref}\PY{p}{,} \PY{n}{sigma\PYZus{}ref} \PY{o}{=} \PY{l+m+mf}{0.0}\PY{p}{,} \PY{l+m+mf}{0.3}\PY{p}{,} \PY{l+m+mf}{1.0}

\PY{c+c1}{\PYZsh{} Parâmetros variados em cada subplot}
\PY{n}{variations} \PY{o}{=} \PY{p}{\PYZob{}}
    \PY{l+s+s2}{\PYZdq{}}\PY{l+s+s2}{u}\PY{l+s+s2}{\PYZdq{}}\PY{p}{:} \PY{p}{\PYZob{}}
        \PY{l+s+s2}{\PYZdq{}}\PY{l+s+s2}{values}\PY{l+s+s2}{\PYZdq{}}\PY{p}{:} \PY{n}{np}\PY{o}{.}\PY{n}{linspace}\PY{p}{(}\PY{l+m+mf}{0.5}\PY{p}{,} \PY{l+m+mf}{4.5}\PY{p}{,} \PY{n}{n\PYZus{}curves}\PY{p}{)}\PY{p}{,}
        \PY{l+s+s2}{\PYZdq{}}\PY{l+s+s2}{label}\PY{l+s+s2}{\PYZdq{}}\PY{p}{:} \PY{l+s+s2}{\PYZdq{}}\PY{l+s+s2}{Ponto de corte \PYZdl{}u\PYZdl{} (localização)}\PY{l+s+s2}{\PYZdq{}}\PY{p}{,}
    \PY{p}{\PYZcb{}}\PY{p}{,}
    \PY{l+s+s2}{\PYZdq{}}\PY{l+s+s2}{k}\PY{l+s+s2}{\PYZdq{}}\PY{p}{:} \PY{p}{\PYZob{}}
        \PY{l+s+s2}{\PYZdq{}}\PY{l+s+s2}{values}\PY{l+s+s2}{\PYZdq{}}\PY{p}{:} \PY{n}{np}\PY{o}{.}\PY{n}{linspace}\PY{p}{(}\PY{o}{\PYZhy{}}\PY{l+m+mf}{1.0}\PY{p}{,} \PY{l+m+mf}{1.0}\PY{p}{,} \PY{n}{n\PYZus{}curves}\PY{p}{)}\PY{p}{,}
        \PY{l+s+s2}{\PYZdq{}}\PY{l+s+s2}{label}\PY{l+s+s2}{\PYZdq{}}\PY{p}{:} \PY{l+s+s2}{\PYZdq{}}\PY{l+s+s2}{Forma \PYZdl{}k\PYZdl{}}\PY{l+s+s2}{\PYZdq{}}\PY{p}{,}
    \PY{p}{\PYZcb{}}\PY{p}{,}
    \PY{l+s+s2}{\PYZdq{}}\PY{l+s+s2}{sigma}\PY{l+s+s2}{\PYZdq{}}\PY{p}{:} \PY{p}{\PYZob{}}
        \PY{l+s+s2}{\PYZdq{}}\PY{l+s+s2}{values}\PY{l+s+s2}{\PYZdq{}}\PY{p}{:} \PY{n}{np}\PY{o}{.}\PY{n}{linspace}\PY{p}{(}\PY{l+m+mf}{0.5}\PY{p}{,} \PY{l+m+mf}{3.0}\PY{p}{,} \PY{n}{n\PYZus{}curves}\PY{p}{)}\PY{p}{,}
        \PY{l+s+s2}{\PYZdq{}}\PY{l+s+s2}{label}\PY{l+s+s2}{\PYZdq{}}\PY{p}{:} \PY{l+s+sa}{r}\PY{l+s+s2}{\PYZdq{}}\PY{l+s+s2}{Escala \PYZdl{}}\PY{l+s+s2}{\PYZbs{}}\PY{l+s+s2}{sigma\PYZdl{}}\PY{l+s+s2}{\PYZdq{}}\PY{p}{,}
    \PY{p}{\PYZcb{}}\PY{p}{,}
\PY{p}{\PYZcb{}}

\PY{n}{cmap} \PY{o}{=} \PY{n}{plt}\PY{o}{.}\PY{n}{get\PYZus{}cmap}\PY{p}{(}\PY{l+s+s2}{\PYZdq{}}\PY{l+s+s2}{magma}\PY{l+s+s2}{\PYZdq{}}\PY{p}{)}
\PY{n}{colors} \PY{o}{=} \PY{p}{[}
    \PY{n}{to\PYZus{}hex}\PY{p}{(}\PY{n}{c}\PY{p}{)} \PY{k}{for} \PY{n}{c} \PY{o+ow}{in} \PY{n}{cmap}\PY{p}{(}\PY{n}{np}\PY{o}{.}\PY{n}{linspace}\PY{p}{(}\PY{l+m+mf}{0.10}\PY{p}{,} \PY{l+m+mf}{0.85}\PY{p}{,} \PY{n}{n\PYZus{}curves}\PY{p}{)}\PY{p}{)}
\PY{p}{]}


\PY{c+c1}{\PYZsh{} \PYZhy{}\PYZhy{}\PYZhy{}\PYZhy{}\PYZhy{}\PYZhy{}\PYZhy{}\PYZhy{}\PYZhy{}\PYZhy{}\PYZhy{}\PYZhy{}\PYZhy{}\PYZhy{}\PYZhy{}\PYZhy{}\PYZhy{}\PYZhy{}\PYZhy{}\PYZhy{}\PYZhy{}\PYZhy{}\PYZhy{}\PYZhy{}\PYZhy{}\PYZhy{}\PYZhy{}\PYZhy{}\PYZhy{}\PYZhy{}\PYZhy{}\PYZhy{}\PYZhy{}\PYZhy{}\PYZhy{}\PYZhy{}\PYZhy{}\PYZhy{}\PYZhy{}\PYZhy{}\PYZhy{}\PYZhy{}\PYZhy{}\PYZhy{}\PYZhy{}\PYZhy{}\PYZhy{}\PYZhy{}\PYZhy{}\PYZhy{}\PYZhy{}\PYZhy{}\PYZhy{}\PYZhy{}\PYZhy{}\PYZhy{}\PYZhy{}\PYZhy{}\PYZhy{}\PYZhy{}\PYZhy{}\PYZhy{}\PYZhy{}\PYZhy{}\PYZhy{}\PYZhy{}\PYZhy{}\PYZhy{}\PYZhy{}\PYZhy{}\PYZhy{}\PYZhy{}\PYZhy{}}
\PY{c+c1}{\PYZsh{} Resumo da configuração}
\PY{c+c1}{\PYZsh{} \PYZhy{}\PYZhy{}\PYZhy{}\PYZhy{}\PYZhy{}\PYZhy{}\PYZhy{}\PYZhy{}\PYZhy{}\PYZhy{}\PYZhy{}\PYZhy{}\PYZhy{}\PYZhy{}\PYZhy{}\PYZhy{}\PYZhy{}\PYZhy{}\PYZhy{}\PYZhy{}\PYZhy{}\PYZhy{}\PYZhy{}\PYZhy{}\PYZhy{}\PYZhy{}\PYZhy{}\PYZhy{}\PYZhy{}\PYZhy{}\PYZhy{}\PYZhy{}\PYZhy{}\PYZhy{}\PYZhy{}\PYZhy{}\PYZhy{}\PYZhy{}\PYZhy{}\PYZhy{}\PYZhy{}\PYZhy{}\PYZhy{}\PYZhy{}\PYZhy{}\PYZhy{}\PYZhy{}\PYZhy{}\PYZhy{}\PYZhy{}\PYZhy{}\PYZhy{}\PYZhy{}\PYZhy{}\PYZhy{}\PYZhy{}\PYZhy{}\PYZhy{}\PYZhy{}\PYZhy{}\PYZhy{}\PYZhy{}\PYZhy{}\PYZhy{}\PYZhy{}\PYZhy{}\PYZhy{}\PYZhy{}\PYZhy{}\PYZhy{}\PYZhy{}\PYZhy{}\PYZhy{}}

\PY{n+nb}{print}\PY{p}{(}\PY{l+s+s2}{\PYZdq{}}\PY{l+s+s2}{\PYZsh{}}\PY{l+s+s2}{\PYZdq{}} \PY{o}{*} \PY{l+m+mi}{100}\PY{p}{)}
\PY{n+nb}{print}\PY{p}{(}\PY{l+s+s2}{\PYZdq{}}\PY{l+s+s2}{DISTRIBUIÇÃO PARETO GENERALIZADA: f(x) = (1/sigma) [1 + k (x\PYZhy{}u)/sigma]\PYZca{}(\PYZhy{}1/k \PYZhy{} 1)}\PY{l+s+s2}{\PYZdq{}}\PY{p}{)}
\PY{n+nb}{print}\PY{p}{(}\PY{l+s+s2}{\PYZdq{}}\PY{l+s+s2}{\PYZsh{}}\PY{l+s+s2}{\PYZdq{}} \PY{o}{*} \PY{l+m+mi}{100}\PY{p}{)}
\PY{n+nb}{print}\PY{p}{(}\PY{l+s+sa}{f}\PY{l+s+s2}{\PYZdq{}}\PY{l+s+s2}{Valores de referência: u = }\PY{l+s+si}{\PYZob{}}\PY{n}{u\PYZus{}ref}\PY{l+s+si}{\PYZcb{}}\PY{l+s+s2}{, k = }\PY{l+s+si}{\PYZob{}}\PY{n}{k\PYZus{}ref}\PY{l+s+si}{\PYZcb{}}\PY{l+s+s2}{, sigma = }\PY{l+s+si}{\PYZob{}}\PY{n}{sigma\PYZus{}ref}\PY{l+s+si}{\PYZcb{}}\PY{l+s+s2}{\PYZdq{}}\PY{p}{)}
\PY{n+nb}{print}\PY{p}{(}\PY{l+s+s2}{\PYZdq{}}\PY{l+s+s2}{Suporte: x \PYZgt{}= u (k \PYZgt{}= 0)  |  u \PYZlt{}= x \PYZlt{}= u \PYZhy{} sigma/k (k \PYZlt{} 0)}\PY{l+s+s2}{\PYZdq{}}\PY{p}{)}
\PY{k}{for} \PY{n}{name}\PY{p}{,} \PY{n}{v} \PY{o+ow}{in} \PY{n}{variations}\PY{o}{.}\PY{n}{items}\PY{p}{(}\PY{p}{)}\PY{p}{:}
    \PY{n+nb}{print}\PY{p}{(}\PY{l+s+sa}{f}\PY{l+s+s2}{\PYZdq{}}\PY{l+s+si}{\PYZob{}}\PY{n}{name}\PY{l+s+si}{:}\PY{l+s+s2}{\PYZlt{}6}\PY{l+s+si}{\PYZcb{}}\PY{l+s+s2}{ varia de }\PY{l+s+si}{\PYZob{}}\PY{n}{v}\PY{p}{[}\PY{l+s+s1}{\PYZsq{}}\PY{l+s+s1}{values}\PY{l+s+s1}{\PYZsq{}}\PY{p}{]}\PY{p}{[}\PY{l+m+mi}{0}\PY{p}{]}\PY{l+s+si}{:}\PY{l+s+s2}{.2f}\PY{l+s+si}{\PYZcb{}}\PY{l+s+s2}{ a }\PY{l+s+si}{\PYZob{}}\PY{n}{v}\PY{p}{[}\PY{l+s+s1}{\PYZsq{}}\PY{l+s+s1}{values}\PY{l+s+s1}{\PYZsq{}}\PY{p}{]}\PY{p}{[}\PY{o}{\PYZhy{}}\PY{l+m+mi}{1}\PY{p}{]}\PY{l+s+si}{:}\PY{l+s+s2}{.2f}\PY{l+s+si}{\PYZcb{}}\PY{l+s+s2}{ }\PY{l+s+s2}{\PYZdq{}}\PY{p}{)}
    \PY{n+nb}{print}\PY{p}{(}\PY{l+s+sa}{f}\PY{l+s+s2}{\PYZdq{}}\PY{l+s+s2}{(}\PY{l+s+si}{\PYZob{}}\PY{n}{n\PYZus{}curves}\PY{l+s+si}{\PYZcb{}}\PY{l+s+s2}{ curvas)}\PY{l+s+s2}{\PYZdq{}}\PY{p}{)}
\PY{n+nb}{print}\PY{p}{(}\PY{l+s+s2}{\PYZdq{}}\PY{l+s+s2}{\PYZsh{}}\PY{l+s+s2}{\PYZdq{}} \PY{o}{*} \PY{l+m+mi}{100}\PY{p}{)}


\PY{c+c1}{\PYZsh{} \PYZhy{}\PYZhy{}\PYZhy{}\PYZhy{}\PYZhy{}\PYZhy{}\PYZhy{}\PYZhy{}\PYZhy{}\PYZhy{}\PYZhy{}\PYZhy{}\PYZhy{}\PYZhy{}\PYZhy{}\PYZhy{}\PYZhy{}\PYZhy{}\PYZhy{}\PYZhy{}\PYZhy{}\PYZhy{}\PYZhy{}\PYZhy{}\PYZhy{}\PYZhy{}\PYZhy{}\PYZhy{}\PYZhy{}\PYZhy{}\PYZhy{}\PYZhy{}\PYZhy{}\PYZhy{}\PYZhy{}\PYZhy{}\PYZhy{}\PYZhy{}\PYZhy{}\PYZhy{}\PYZhy{}\PYZhy{}\PYZhy{}\PYZhy{}\PYZhy{}\PYZhy{}\PYZhy{}\PYZhy{}\PYZhy{}\PYZhy{}\PYZhy{}\PYZhy{}\PYZhy{}\PYZhy{}\PYZhy{}\PYZhy{}\PYZhy{}\PYZhy{}\PYZhy{}\PYZhy{}\PYZhy{}\PYZhy{}\PYZhy{}\PYZhy{}\PYZhy{}\PYZhy{}\PYZhy{}\PYZhy{}\PYZhy{}\PYZhy{}\PYZhy{}\PYZhy{}\PYZhy{}}
\PY{c+c1}{\PYZsh{} Visualização}
\PY{c+c1}{\PYZsh{} \PYZhy{}\PYZhy{}\PYZhy{}\PYZhy{}\PYZhy{}\PYZhy{}\PYZhy{}\PYZhy{}\PYZhy{}\PYZhy{}\PYZhy{}\PYZhy{}\PYZhy{}\PYZhy{}\PYZhy{}\PYZhy{}\PYZhy{}\PYZhy{}\PYZhy{}\PYZhy{}\PYZhy{}\PYZhy{}\PYZhy{}\PYZhy{}\PYZhy{}\PYZhy{}\PYZhy{}\PYZhy{}\PYZhy{}\PYZhy{}\PYZhy{}\PYZhy{}\PYZhy{}\PYZhy{}\PYZhy{}\PYZhy{}\PYZhy{}\PYZhy{}\PYZhy{}\PYZhy{}\PYZhy{}\PYZhy{}\PYZhy{}\PYZhy{}\PYZhy{}\PYZhy{}\PYZhy{}\PYZhy{}\PYZhy{}\PYZhy{}\PYZhy{}\PYZhy{}\PYZhy{}\PYZhy{}\PYZhy{}\PYZhy{}\PYZhy{}\PYZhy{}\PYZhy{}\PYZhy{}\PYZhy{}\PYZhy{}\PYZhy{}\PYZhy{}\PYZhy{}\PYZhy{}\PYZhy{}\PYZhy{}\PYZhy{}\PYZhy{}\PYZhy{}\PYZhy{}\PYZhy{}}

\PY{n}{fig}\PY{p}{,} \PY{n}{axes} \PY{o}{=} \PY{n}{plt}\PY{o}{.}\PY{n}{subplots}\PY{p}{(}
    \PY{n}{nrows}\PY{o}{=}\PY{l+m+mi}{3}\PY{p}{,} \PY{n}{ncols}\PY{o}{=}\PY{l+m+mi}{1}\PY{p}{,}
    \PY{n}{figsize}\PY{o}{=}\PY{p}{(}\PY{l+m+mi}{18}\PY{p}{,} \PY{l+m+mi}{16}\PY{p}{)}
\PY{p}{)}

\PY{n}{fig}\PY{o}{.}\PY{n}{suptitle}\PY{p}{(}
    \PY{l+s+s2}{\PYZdq{}}\PY{l+s+s2}{Distribuição Pareto generalizada: efeito de cada parâmetro}\PY{l+s+s2}{\PYZdq{}}\PY{p}{,}
    \PY{n}{fontsize}\PY{o}{=}\PY{l+m+mi}{22}\PY{p}{,}
    \PY{n}{y}\PY{o}{=}\PY{l+m+mf}{0.995}
\PY{p}{)}

\PY{k}{for} \PY{n}{ax}\PY{p}{,} \PY{p}{(}\PY{n}{name}\PY{p}{,} \PY{n}{var}\PY{p}{)} \PY{o+ow}{in} \PY{n+nb}{zip}\PY{p}{(}\PY{n}{axes}\PY{p}{,} \PY{n}{variations}\PY{o}{.}\PY{n}{items}\PY{p}{(}\PY{p}{)}\PY{p}{)}\PY{p}{:}

    \PY{k}{for} \PY{n}{value}\PY{p}{,} \PY{n}{color} \PY{o+ow}{in} \PY{n+nb}{zip}\PY{p}{(}\PY{n}{var}\PY{p}{[}\PY{l+s+s2}{\PYZdq{}}\PY{l+s+s2}{values}\PY{l+s+s2}{\PYZdq{}}\PY{p}{]}\PY{p}{,} \PY{n}{colors}\PY{p}{)}\PY{p}{:}

        \PY{n}{u}\PY{p}{,} \PY{n}{k}\PY{p}{,} \PY{n}{sigma} \PY{o}{=} \PY{n}{u\PYZus{}ref}\PY{p}{,} \PY{n}{k\PYZus{}ref}\PY{p}{,} \PY{n}{sigma\PYZus{}ref}
        \PY{k}{if} \PY{n}{name} \PY{o}{==} \PY{l+s+s2}{\PYZdq{}}\PY{l+s+s2}{u}\PY{l+s+s2}{\PYZdq{}}\PY{p}{:}
            \PY{n}{u} \PY{o}{=} \PY{n}{value}
        \PY{k}{elif} \PY{n}{name} \PY{o}{==} \PY{l+s+s2}{\PYZdq{}}\PY{l+s+s2}{k}\PY{l+s+s2}{\PYZdq{}}\PY{p}{:}
            \PY{n}{k} \PY{o}{=} \PY{n}{value}
        \PY{k}{else}\PY{p}{:}
            \PY{n}{sigma} \PY{o}{=} \PY{n}{value}

        \PY{c+c1}{\PYZsh{} scipy: c = k (forma), loc = u (corte), scale = sigma (escala)}
        \PY{n}{pdf} \PY{o}{=} \PY{n}{stats}\PY{o}{.}\PY{n}{genpareto}\PY{o}{.}\PY{n}{pdf}\PY{p}{(}\PY{n}{x}\PY{p}{,} \PY{n}{c}\PY{o}{=}\PY{n}{k}\PY{p}{,} \PY{n}{loc}\PY{o}{=}\PY{n}{u}\PY{p}{,} \PY{n}{scale}\PY{o}{=}\PY{n}{sigma}\PY{p}{)}

        \PY{n}{ax}\PY{o}{.}\PY{n}{plot}\PY{p}{(}
            \PY{n}{x}\PY{p}{,} \PY{n}{pdf}\PY{p}{,}
            \PY{n}{color}\PY{o}{=}\PY{n}{color}\PY{p}{,} \PY{n}{linewidth}\PY{o}{=}\PY{l+m+mf}{2.5}\PY{p}{,}
            \PY{n}{label}\PY{o}{=}\PY{l+s+sa}{f}\PY{l+s+s2}{\PYZdq{}}\PY{l+s+si}{\PYZob{}}\PY{n}{name}\PY{l+s+si}{\PYZcb{}}\PY{l+s+s2}{ = }\PY{l+s+si}{\PYZob{}}\PY{n}{value}\PY{l+s+si}{:}\PY{l+s+s2}{.2f}\PY{l+s+si}{\PYZcb{}}\PY{l+s+s2}{\PYZdq{}}
        \PY{p}{)}

    \PY{n}{fixed} \PY{o}{=} \PY{p}{\PYZob{}}\PY{l+s+s2}{\PYZdq{}}\PY{l+s+s2}{u}\PY{l+s+s2}{\PYZdq{}}\PY{p}{:} \PY{n}{u\PYZus{}ref}\PY{p}{,} \PY{l+s+s2}{\PYZdq{}}\PY{l+s+s2}{k}\PY{l+s+s2}{\PYZdq{}}\PY{p}{:} \PY{n}{k\PYZus{}ref}\PY{p}{,} \PY{l+s+s2}{\PYZdq{}}\PY{l+s+s2}{sigma}\PY{l+s+s2}{\PYZdq{}}\PY{p}{:} \PY{n}{sigma\PYZus{}ref}\PY{p}{\PYZcb{}}
    \PY{n}{fixed}\PY{o}{.}\PY{n}{pop}\PY{p}{(}\PY{n}{name}\PY{p}{)}
    \PY{n}{fixed\PYZus{}txt} \PY{o}{=} \PY{l+s+s2}{\PYZdq{}}\PY{l+s+s2}{, }\PY{l+s+s2}{\PYZdq{}}\PY{o}{.}\PY{n}{join}\PY{p}{(}\PY{l+s+sa}{f}\PY{l+s+s2}{\PYZdq{}}\PY{l+s+si}{\PYZob{}}\PY{n}{p}\PY{l+s+si}{\PYZcb{}}\PY{l+s+s2}{ = }\PY{l+s+si}{\PYZob{}}\PY{n}{v}\PY{l+s+si}{\PYZcb{}}\PY{l+s+s2}{\PYZdq{}} \PY{k}{for} \PY{n}{p}\PY{p}{,} \PY{n}{v} \PY{o+ow}{in} \PY{n}{fixed}\PY{o}{.}\PY{n}{items}\PY{p}{(}\PY{p}{)}\PY{p}{)}

    \PY{n}{ax}\PY{o}{.}\PY{n}{set\PYZus{}xlabel}\PY{p}{(}\PY{l+s+s2}{\PYZdq{}}\PY{l+s+s2}{x}\PY{l+s+s2}{\PYZdq{}}\PY{p}{,} \PY{n}{fontsize}\PY{o}{=}\PY{l+m+mi}{15}\PY{p}{)}
    \PY{n}{ax}\PY{o}{.}\PY{n}{set\PYZus{}ylabel}\PY{p}{(}\PY{l+s+s2}{\PYZdq{}}\PY{l+s+s2}{densidade}\PY{l+s+s2}{\PYZdq{}}\PY{p}{,} \PY{n}{fontsize}\PY{o}{=}\PY{l+m+mi}{15}\PY{p}{)}
    \PY{n}{ax}\PY{o}{.}\PY{n}{set\PYZus{}title}\PY{p}{(}
        \PY{l+s+sa}{f}\PY{l+s+s2}{\PYZdq{}}\PY{l+s+s2}{Variando }\PY{l+s+si}{\PYZob{}}\PY{n}{var}\PY{p}{[}\PY{l+s+s1}{\PYZsq{}}\PY{l+s+s1}{label}\PY{l+s+s1}{\PYZsq{}}\PY{p}{]}\PY{l+s+si}{\PYZcb{}}\PY{l+s+s2}{  (fixos: }\PY{l+s+si}{\PYZob{}}\PY{n}{fixed\PYZus{}txt}\PY{l+s+si}{\PYZcb{}}\PY{l+s+s2}{)}\PY{l+s+s2}{\PYZdq{}}\PY{p}{,}
        \PY{n}{fontsize}\PY{o}{=}\PY{l+m+mi}{17}
    \PY{p}{)}
    \PY{n}{ax}\PY{o}{.}\PY{n}{set\PYZus{}xlim}\PY{p}{(}\PY{n}{x}\PY{p}{[}\PY{l+m+mi}{0}\PY{p}{]}\PY{p}{,} \PY{n}{x}\PY{p}{[}\PY{o}{\PYZhy{}}\PY{l+m+mi}{1}\PY{p}{]}\PY{p}{)}
    \PY{n}{ax}\PY{o}{.}\PY{n}{tick\PYZus{}params}\PY{p}{(}\PY{n}{axis}\PY{o}{=}\PY{l+s+s2}{\PYZdq{}}\PY{l+s+s2}{both}\PY{l+s+s2}{\PYZdq{}}\PY{p}{,} \PY{n}{labelsize}\PY{o}{=}\PY{l+m+mi}{13}\PY{p}{)}
    \PY{n}{ax}\PY{o}{.}\PY{n}{grid}\PY{p}{(}\PY{n}{alpha}\PY{o}{=}\PY{l+m+mf}{0.25}\PY{p}{,} \PY{n}{linestyle}\PY{o}{=}\PY{l+s+s2}{\PYZdq{}}\PY{l+s+s2}{\PYZhy{}\PYZhy{}}\PY{l+s+s2}{\PYZdq{}}\PY{p}{)}
    \PY{n}{ax}\PY{o}{.}\PY{n}{legend}\PY{p}{(}\PY{n}{fontsize}\PY{o}{=}\PY{l+m+mi}{12}\PY{p}{,} \PY{n}{loc}\PY{o}{=}\PY{l+s+s2}{\PYZdq{}}\PY{l+s+s2}{center left}\PY{l+s+s2}{\PYZdq{}}\PY{p}{,} \PY{n}{bbox\PYZus{}to\PYZus{}anchor}\PY{o}{=}\PY{p}{(}\PY{l+m+mf}{1.01}\PY{p}{,} \PY{l+m+mf}{0.5}\PY{p}{)}\PY{p}{)}

\PY{n}{plt}\PY{o}{.}\PY{n}{tight\PYZus{}layout}\PY{p}{(}\PY{n}{rect}\PY{o}{=}\PY{p}{[}\PY{l+m+mi}{0}\PY{p}{,} \PY{l+m+mi}{0}\PY{p}{,} \PY{l+m+mi}{1}\PY{p}{,} \PY{l+m+mf}{0.98}\PY{p}{]}\PY{p}{)}
\PY{n}{plt}\PY{o}{.}\PY{n}{show}\PY{p}{(}\PY{p}{)}
\end{Verbatim}
\end{tcolorbox}

    \begin{Verbatim}[commandchars=\\\{\}]
\#\#\#\#\#\#\#\#\#\#\#\#\#\#\#\#\#\#\#\#\#\#\#\#\#\#\#\#\#\#\#\#\#\#\#\#\#\#\#\#\#\#\#\#\#\#\#\#\#\#\#\#\#\#\#\#\#\#\#\#\#\#\#\#\#\#\#\#\#\#\#\#\#\#\#\#\#\#\#\#
\#\#\#\#\#\#\#\#\#\#\#\#\#\#\#\#\#\#\#\#
DISTRIBUIÇÃO PARETO GENERALIZADA: f(x) = (1/sigma) [1 + k (x-u)/sigma]\^{}(-1/k -
1)
\#\#\#\#\#\#\#\#\#\#\#\#\#\#\#\#\#\#\#\#\#\#\#\#\#\#\#\#\#\#\#\#\#\#\#\#\#\#\#\#\#\#\#\#\#\#\#\#\#\#\#\#\#\#\#\#\#\#\#\#\#\#\#\#\#\#\#\#\#\#\#\#\#\#\#\#\#\#\#\#
\#\#\#\#\#\#\#\#\#\#\#\#\#\#\#\#\#\#\#\#
Valores de referência: u = 0.0, k = 0.3, sigma = 1.0
Suporte: x >= u (k >= 0)  |  u <= x <= u - sigma/k (k < 0)
u      varia de 0.50 a 4.50
(10 curvas)
k      varia de -1.00 a 1.00
(10 curvas)
sigma  varia de 0.50 a 3.00
(10 curvas)
\#\#\#\#\#\#\#\#\#\#\#\#\#\#\#\#\#\#\#\#\#\#\#\#\#\#\#\#\#\#\#\#\#\#\#\#\#\#\#\#\#\#\#\#\#\#\#\#\#\#\#\#\#\#\#\#\#\#\#\#\#\#\#\#\#\#\#\#\#\#\#\#\#\#\#\#\#\#\#\#
\#\#\#\#\#\#\#\#\#\#\#\#\#\#\#\#\#\#\#\#
    \end{Verbatim}

    \begin{center}
    \adjustimage{max size={0.9\linewidth}{0.9\paperheight}}{notebook_corrigido_files/notebook_corrigido_4_1.png}
    \end{center}
    { \hspace*{\fill} \\}
    
    \begin{tcolorbox}[breakable, size=fbox, boxrule=1pt, pad at break*=1mm,colback=cellbackground, colframe=cellborder]
\prompt{In}{incolor}{28}{\boxspacing}
\begin{Verbatim}[commandchars=\\\{\}]
\PY{c+c1}{\PYZsh{} \PYZhy{}\PYZhy{}\PYZhy{}\PYZhy{}\PYZhy{}\PYZhy{}\PYZhy{}\PYZhy{}\PYZhy{}\PYZhy{}\PYZhy{}\PYZhy{}\PYZhy{}\PYZhy{}\PYZhy{}\PYZhy{}\PYZhy{}\PYZhy{}\PYZhy{}\PYZhy{}\PYZhy{}\PYZhy{}\PYZhy{}\PYZhy{}\PYZhy{}\PYZhy{}\PYZhy{}\PYZhy{}\PYZhy{}\PYZhy{}\PYZhy{}\PYZhy{}\PYZhy{}\PYZhy{}\PYZhy{}\PYZhy{}\PYZhy{}\PYZhy{}\PYZhy{}\PYZhy{}\PYZhy{}\PYZhy{}\PYZhy{}\PYZhy{}\PYZhy{}\PYZhy{}\PYZhy{}\PYZhy{}\PYZhy{}\PYZhy{}\PYZhy{}\PYZhy{}\PYZhy{}\PYZhy{}\PYZhy{}\PYZhy{}\PYZhy{}\PYZhy{}\PYZhy{}\PYZhy{}\PYZhy{}\PYZhy{}\PYZhy{}\PYZhy{}\PYZhy{}\PYZhy{}\PYZhy{}\PYZhy{}\PYZhy{}\PYZhy{}\PYZhy{}\PYZhy{}\PYZhy{}}
\PY{c+c1}{\PYZsh{} Configurações}
\PY{c+c1}{\PYZsh{} \PYZhy{}\PYZhy{}\PYZhy{}\PYZhy{}\PYZhy{}\PYZhy{}\PYZhy{}\PYZhy{}\PYZhy{}\PYZhy{}\PYZhy{}\PYZhy{}\PYZhy{}\PYZhy{}\PYZhy{}\PYZhy{}\PYZhy{}\PYZhy{}\PYZhy{}\PYZhy{}\PYZhy{}\PYZhy{}\PYZhy{}\PYZhy{}\PYZhy{}\PYZhy{}\PYZhy{}\PYZhy{}\PYZhy{}\PYZhy{}\PYZhy{}\PYZhy{}\PYZhy{}\PYZhy{}\PYZhy{}\PYZhy{}\PYZhy{}\PYZhy{}\PYZhy{}\PYZhy{}\PYZhy{}\PYZhy{}\PYZhy{}\PYZhy{}\PYZhy{}\PYZhy{}\PYZhy{}\PYZhy{}\PYZhy{}\PYZhy{}\PYZhy{}\PYZhy{}\PYZhy{}\PYZhy{}\PYZhy{}\PYZhy{}\PYZhy{}\PYZhy{}\PYZhy{}\PYZhy{}\PYZhy{}\PYZhy{}\PYZhy{}\PYZhy{}\PYZhy{}\PYZhy{}\PYZhy{}\PYZhy{}\PYZhy{}\PYZhy{}\PYZhy{}\PYZhy{}\PYZhy{}}

\PY{n}{rng} \PY{o}{=} \PY{n}{np}\PY{o}{.}\PY{n}{random}\PY{o}{.}\PY{n}{default\PYZus{}rng}\PY{p}{(}\PY{l+m+mi}{121}\PY{p}{)}

\PY{n}{n\PYZus{}samples} \PY{o}{=} \PY{l+m+mi}{5000}
\PY{n}{reff} \PY{o}{=} \PY{l+m+mf}{1.0}
\PY{n}{log\PYZus{}y} \PY{o}{=} \PY{k+kc}{True}   \PY{c+c1}{\PYZsh{} escala logarítmica no eixo y dos gráficos de pesos}

\PY{c+c1}{\PYZsh{} Forma funcional única: log w = k * E, com E \PYZti{} Exponential(1).}
\PY{c+c1}{\PYZsh{} Escalar os log\PYZhy{}pesos por k equivale a log w \PYZti{} Exponential(taxa = 1/k), logo}
\PY{c+c1}{\PYZsh{}   P(w \PYZgt{} t) = t\PYZca{}(\PYZhy{}1/k)  \PYZhy{}\PYZgt{}  w \PYZti{} Pareto, cauda de Pareto generalizada com forma k.}
\PY{c+c1}{\PYZsh{}   k maior (escala maior) \PYZhy{}\PYZgt{} decaimento mais lento  \PYZhy{}\PYZgt{} cauda mais pesada}
\PY{c+c1}{\PYZsh{}   k menor (escala menor) \PYZhy{}\PYZgt{} decaimento mais rápido \PYZhy{}\PYZgt{} cauda mais leve}
\PY{n}{k\PYZus{}values} \PY{o}{=} \PY{p}{[}\PY{o}{\PYZhy{}}\PY{l+m+mf}{0.25}\PY{p}{,} \PY{l+m+mf}{0.05}\PY{p}{,} \PY{l+m+mf}{0.5}\PY{p}{,} \PY{l+m+mf}{0.6}\PY{p}{,} \PY{l+m+mf}{1.0}\PY{p}{]}

\PY{n}{scenarios} \PY{o}{=} \PY{p}{\PYZob{}}
    \PY{l+s+sa}{f}\PY{l+s+s2}{\PYZdq{}}\PY{l+s+s2}{Cenário }\PY{l+s+si}{\PYZob{}}\PY{n}{i}\PY{l+s+si}{\PYZcb{}}\PY{l+s+s2}{\PYZdq{}}\PY{p}{:} \PY{p}{\PYZob{}}\PY{l+s+s2}{\PYZdq{}}\PY{l+s+s2}{k}\PY{l+s+s2}{\PYZdq{}}\PY{p}{:} \PY{n}{k}\PY{p}{\PYZcb{}}
    \PY{k}{for} \PY{n}{i}\PY{p}{,} \PY{n}{k} \PY{o+ow}{in} \PY{n+nb}{enumerate}\PY{p}{(}\PY{n}{k\PYZus{}values}\PY{p}{,} \PY{n}{start}\PY{o}{=}\PY{l+m+mi}{1}\PY{p}{)}
\PY{p}{\PYZcb{}}

\PY{n}{cmap} \PY{o}{=} \PY{n}{plt}\PY{o}{.}\PY{n}{get\PYZus{}cmap}\PY{p}{(}\PY{l+s+s2}{\PYZdq{}}\PY{l+s+s2}{magma}\PY{l+s+s2}{\PYZdq{}}\PY{p}{)}
\PY{n}{colors} \PY{o}{=} \PY{p}{[}
    \PY{n}{to\PYZus{}hex}\PY{p}{(}\PY{n}{c}\PY{p}{)} \PY{k}{for} \PY{n}{c} \PY{o+ow}{in} \PY{n}{cmap}\PY{p}{(}\PY{n}{np}\PY{o}{.}\PY{n}{linspace}\PY{p}{(}\PY{l+m+mf}{0.15}\PY{p}{,} \PY{l+m+mf}{0.80}\PY{p}{,} \PY{n+nb}{len}\PY{p}{(}\PY{n}{scenarios}\PY{p}{)}\PY{p}{)}\PY{p}{)}
\PY{p}{]}


\PY{c+c1}{\PYZsh{} \PYZhy{}\PYZhy{}\PYZhy{}\PYZhy{}\PYZhy{}\PYZhy{}\PYZhy{}\PYZhy{}\PYZhy{}\PYZhy{}\PYZhy{}\PYZhy{}\PYZhy{}\PYZhy{}\PYZhy{}\PYZhy{}\PYZhy{}\PYZhy{}\PYZhy{}\PYZhy{}\PYZhy{}\PYZhy{}\PYZhy{}\PYZhy{}\PYZhy{}\PYZhy{}\PYZhy{}\PYZhy{}\PYZhy{}\PYZhy{}\PYZhy{}\PYZhy{}\PYZhy{}\PYZhy{}\PYZhy{}\PYZhy{}\PYZhy{}\PYZhy{}\PYZhy{}\PYZhy{}\PYZhy{}\PYZhy{}\PYZhy{}\PYZhy{}\PYZhy{}\PYZhy{}\PYZhy{}\PYZhy{}\PYZhy{}\PYZhy{}\PYZhy{}\PYZhy{}\PYZhy{}\PYZhy{}\PYZhy{}\PYZhy{}\PYZhy{}\PYZhy{}\PYZhy{}\PYZhy{}\PYZhy{}\PYZhy{}\PYZhy{}\PYZhy{}\PYZhy{}\PYZhy{}\PYZhy{}\PYZhy{}\PYZhy{}\PYZhy{}\PYZhy{}\PYZhy{}\PYZhy{}}
\PY{c+c1}{\PYZsh{} Funções auxiliares}
\PY{c+c1}{\PYZsh{} \PYZhy{}\PYZhy{}\PYZhy{}\PYZhy{}\PYZhy{}\PYZhy{}\PYZhy{}\PYZhy{}\PYZhy{}\PYZhy{}\PYZhy{}\PYZhy{}\PYZhy{}\PYZhy{}\PYZhy{}\PYZhy{}\PYZhy{}\PYZhy{}\PYZhy{}\PYZhy{}\PYZhy{}\PYZhy{}\PYZhy{}\PYZhy{}\PYZhy{}\PYZhy{}\PYZhy{}\PYZhy{}\PYZhy{}\PYZhy{}\PYZhy{}\PYZhy{}\PYZhy{}\PYZhy{}\PYZhy{}\PYZhy{}\PYZhy{}\PYZhy{}\PYZhy{}\PYZhy{}\PYZhy{}\PYZhy{}\PYZhy{}\PYZhy{}\PYZhy{}\PYZhy{}\PYZhy{}\PYZhy{}\PYZhy{}\PYZhy{}\PYZhy{}\PYZhy{}\PYZhy{}\PYZhy{}\PYZhy{}\PYZhy{}\PYZhy{}\PYZhy{}\PYZhy{}\PYZhy{}\PYZhy{}\PYZhy{}\PYZhy{}\PYZhy{}\PYZhy{}\PYZhy{}\PYZhy{}\PYZhy{}\PYZhy{}\PYZhy{}\PYZhy{}\PYZhy{}\PYZhy{}}

\PY{k}{def}\PY{+w}{ }\PY{n+nf}{generate\PYZus{}log\PYZus{}weights}\PY{p}{(}\PY{n}{spec}\PY{p}{,} \PY{n}{size}\PY{p}{,} \PY{n}{rng}\PY{p}{)}\PY{p}{:}
\PY{+w}{    }\PY{l+s+sd}{\PYZdq{}\PYZdq{}\PYZdq{}Log\PYZhy{}pesos: exponencial padrão escalada por spec[\PYZsq{}k\PYZsq{}].\PYZdq{}\PYZdq{}\PYZdq{}}
    \PY{k}{return} \PY{n}{spec}\PY{p}{[}\PY{l+s+s2}{\PYZdq{}}\PY{l+s+s2}{k}\PY{l+s+s2}{\PYZdq{}}\PY{p}{]} \PY{o}{*} \PY{n}{rng}\PY{o}{.}\PY{n}{exponential}\PY{p}{(}\PY{n}{scale}\PY{o}{=}\PY{l+m+mf}{1.0}\PY{p}{,} \PY{n}{size}\PY{o}{=}\PY{n}{size}\PY{p}{)}


\PY{k}{def}\PY{+w}{ }\PY{n+nf}{psis\PYZus{}tail\PYZus{}length}\PY{p}{(}\PY{n}{size}\PY{p}{,} \PY{n}{reff}\PY{o}{=}\PY{l+m+mf}{1.0}\PY{p}{)}\PY{p}{:}
\PY{+w}{    }\PY{l+s+sd}{\PYZdq{}\PYZdq{}\PYZdq{}Número de pesos na cauda suavizada pelo PSIS.\PYZdq{}\PYZdq{}\PYZdq{}}
    \PY{k}{return} \PY{n+nb}{int}\PY{p}{(}\PY{n}{np}\PY{o}{.}\PY{n}{ceil}\PY{p}{(}\PY{n+nb}{min}\PY{p}{(}\PY{n}{size} \PY{o}{/} \PY{l+m+mi}{5}\PY{p}{,} \PY{l+m+mi}{3} \PY{o}{*} \PY{n}{np}\PY{o}{.}\PY{n}{sqrt}\PY{p}{(}\PY{n}{size} \PY{o}{/} \PY{n}{reff}\PY{p}{)}\PY{p}{)}\PY{p}{)}\PY{p}{)}


\PY{c+c1}{\PYZsh{} \PYZhy{}\PYZhy{}\PYZhy{}\PYZhy{}\PYZhy{}\PYZhy{}\PYZhy{}\PYZhy{}\PYZhy{}\PYZhy{}\PYZhy{}\PYZhy{}\PYZhy{}\PYZhy{}\PYZhy{}\PYZhy{}\PYZhy{}\PYZhy{}\PYZhy{}\PYZhy{}\PYZhy{}\PYZhy{}\PYZhy{}\PYZhy{}\PYZhy{}\PYZhy{}\PYZhy{}\PYZhy{}\PYZhy{}\PYZhy{}\PYZhy{}\PYZhy{}\PYZhy{}\PYZhy{}\PYZhy{}\PYZhy{}\PYZhy{}\PYZhy{}\PYZhy{}\PYZhy{}\PYZhy{}\PYZhy{}\PYZhy{}\PYZhy{}\PYZhy{}\PYZhy{}\PYZhy{}\PYZhy{}\PYZhy{}\PYZhy{}\PYZhy{}\PYZhy{}\PYZhy{}\PYZhy{}\PYZhy{}\PYZhy{}\PYZhy{}\PYZhy{}\PYZhy{}\PYZhy{}\PYZhy{}\PYZhy{}\PYZhy{}\PYZhy{}\PYZhy{}\PYZhy{}\PYZhy{}\PYZhy{}\PYZhy{}\PYZhy{}\PYZhy{}\PYZhy{}\PYZhy{}}
\PY{c+c1}{\PYZsh{} Resumo da configuração do estudo}
\PY{c+c1}{\PYZsh{} \PYZhy{}\PYZhy{}\PYZhy{}\PYZhy{}\PYZhy{}\PYZhy{}\PYZhy{}\PYZhy{}\PYZhy{}\PYZhy{}\PYZhy{}\PYZhy{}\PYZhy{}\PYZhy{}\PYZhy{}\PYZhy{}\PYZhy{}\PYZhy{}\PYZhy{}\PYZhy{}\PYZhy{}\PYZhy{}\PYZhy{}\PYZhy{}\PYZhy{}\PYZhy{}\PYZhy{}\PYZhy{}\PYZhy{}\PYZhy{}\PYZhy{}\PYZhy{}\PYZhy{}\PYZhy{}\PYZhy{}\PYZhy{}\PYZhy{}\PYZhy{}\PYZhy{}\PYZhy{}\PYZhy{}\PYZhy{}\PYZhy{}\PYZhy{}\PYZhy{}\PYZhy{}\PYZhy{}\PYZhy{}\PYZhy{}\PYZhy{}\PYZhy{}\PYZhy{}\PYZhy{}\PYZhy{}\PYZhy{}\PYZhy{}\PYZhy{}\PYZhy{}\PYZhy{}\PYZhy{}\PYZhy{}\PYZhy{}\PYZhy{}\PYZhy{}\PYZhy{}\PYZhy{}\PYZhy{}\PYZhy{}\PYZhy{}\PYZhy{}\PYZhy{}\PYZhy{}\PYZhy{}}

\PY{n}{tail\PYZus{}length} \PY{o}{=} \PY{n}{psis\PYZus{}tail\PYZus{}length}\PY{p}{(}\PY{n}{n\PYZus{}samples}\PY{p}{,} \PY{n}{reff}\PY{p}{)}

\PY{n+nb}{print}\PY{p}{(}\PY{l+s+s2}{\PYZdq{}}\PY{l+s+s2}{\PYZsh{}}\PY{l+s+s2}{\PYZdq{}} \PY{o}{*} \PY{l+m+mi}{100}\PY{p}{)}
\PY{n+nb}{print}\PY{p}{(}\PY{l+s+s2}{\PYZdq{}}\PY{l+s+s2}{CONFIGURAÇÃO DA SIMULAÇÃO: PSIS SOBRE PESOS DE IMPORTÂNCIA SINTÉTICOS}\PY{l+s+s2}{\PYZdq{}}\PY{p}{)}
\PY{n+nb}{print}\PY{p}{(}\PY{l+s+s2}{\PYZdq{}}\PY{l+s+s2}{\PYZsh{}}\PY{l+s+s2}{\PYZdq{}} \PY{o}{*} \PY{l+m+mi}{100}\PY{p}{)}
\PY{n+nb}{print}\PY{p}{(}\PY{l+s+sa}{f}\PY{l+s+s2}{\PYZdq{}}\PY{l+s+s2}{Amostras:         K = }\PY{l+s+si}{\PYZob{}}\PY{n}{n\PYZus{}samples}\PY{l+s+si}{\PYZcb{}}\PY{l+s+s2}{  (reff = }\PY{l+s+si}{\PYZob{}}\PY{n}{reff}\PY{l+s+si}{\PYZcb{}}\PY{l+s+s2}{)}\PY{l+s+s2}{\PYZdq{}}\PY{p}{)}
\PY{n+nb}{print}\PY{p}{(}\PY{l+s+sa}{f}\PY{l+s+s2}{\PYZdq{}}\PY{l+s+s2}{Cauda do PSIS:    M = }\PY{l+s+si}{\PYZob{}}\PY{n}{tail\PYZus{}length}\PY{l+s+si}{\PYZcb{}}\PY{l+s+s2}{ maiores pesos }\PY{l+s+s2}{\PYZdq{}}
      \PY{l+s+sa}{f}\PY{l+s+s2}{\PYZdq{}}\PY{l+s+s2}{(}\PY{l+s+si}{\PYZob{}}\PY{l+m+mi}{100}\PY{+w}{ }\PY{o}{*}\PY{+w}{ }\PY{n}{tail\PYZus{}length}\PY{+w}{ }\PY{o}{/}\PY{+w}{ }\PY{n}{n\PYZus{}samples}\PY{l+s+si}{:}\PY{l+s+s2}{.1f}\PY{l+s+si}{\PYZcb{}}\PY{l+s+s2}{\PYZpc{} da amostra)}\PY{l+s+s2}{\PYZdq{}}\PY{p}{)}
\PY{n+nb}{print}\PY{p}{(}\PY{l+s+sa}{f}\PY{l+s+s2}{\PYZdq{}}\PY{l+s+s2}{Suavização:       Pareto generalizada ajustada à cauda (arviz.psislw)}\PY{l+s+s2}{\PYZdq{}}\PY{p}{)}
\PY{n+nb}{print}\PY{p}{(}\PY{l+s+s2}{\PYZdq{}}\PY{l+s+s2}{\PYZhy{}}\PY{l+s+s2}{\PYZdq{}} \PY{o}{*} \PY{l+m+mi}{100}\PY{p}{)}
\PY{n+nb}{print}\PY{p}{(}\PY{l+s+s2}{\PYZdq{}}\PY{l+s+s2}{Geração dos pesos brutos (forma funcional única):}\PY{l+s+s2}{\PYZdq{}}\PY{p}{)}
\PY{n+nb}{print}\PY{p}{(}\PY{l+s+s2}{\PYZdq{}}\PY{l+s+s2}{  log w = k * E, com E \PYZti{} Exponencial(1)}\PY{l+s+s2}{\PYZdq{}}\PY{p}{)}
\PY{n+nb}{print}\PY{p}{(}\PY{l+s+s2}{\PYZdq{}}\PY{l+s+s2}{  Escala maior (decaimento mais lento) \PYZhy{}\PYZgt{} cauda mais pesada \PYZhy{}\PYZgt{} k maior.}\PY{l+s+s2}{\PYZdq{}}\PY{p}{)}
\PY{n+nb}{print}\PY{p}{(}\PY{l+s+s2}{\PYZdq{}}\PY{l+s+s2}{  Escala menor (decaimento mais rápido) \PYZhy{}\PYZgt{} cauda mais leve \PYZhy{}\PYZgt{} k menor.}\PY{l+s+s2}{\PYZdq{}}\PY{p}{)}
\PY{n+nb}{print}\PY{p}{(}\PY{l+s+s2}{\PYZdq{}}\PY{l+s+s2}{\PYZsh{}}\PY{l+s+s2}{\PYZdq{}} \PY{o}{*} \PY{l+m+mi}{100}\PY{p}{)}
\PY{n+nb}{print}\PY{p}{(}\PY{l+s+s2}{\PYZdq{}}\PY{l+s+s2}{Leitura: k\PYZhy{}hat \PYZlt{} 0,5 (leve); 0,5\PYZhy{}0,7 (limiar); \PYZgt{} 0,7 (pouco confiável).}\PY{l+s+s2}{\PYZdq{}}\PY{p}{)}
\PY{n+nb}{print}\PY{p}{(}\PY{p}{)}


\PY{c+c1}{\PYZsh{} \PYZhy{}\PYZhy{}\PYZhy{}\PYZhy{}\PYZhy{}\PYZhy{}\PYZhy{}\PYZhy{}\PYZhy{}\PYZhy{}\PYZhy{}\PYZhy{}\PYZhy{}\PYZhy{}\PYZhy{}\PYZhy{}\PYZhy{}\PYZhy{}\PYZhy{}\PYZhy{}\PYZhy{}\PYZhy{}\PYZhy{}\PYZhy{}\PYZhy{}\PYZhy{}\PYZhy{}\PYZhy{}\PYZhy{}\PYZhy{}\PYZhy{}\PYZhy{}\PYZhy{}\PYZhy{}\PYZhy{}\PYZhy{}\PYZhy{}\PYZhy{}\PYZhy{}\PYZhy{}\PYZhy{}\PYZhy{}\PYZhy{}\PYZhy{}\PYZhy{}\PYZhy{}\PYZhy{}\PYZhy{}\PYZhy{}\PYZhy{}\PYZhy{}\PYZhy{}\PYZhy{}\PYZhy{}\PYZhy{}\PYZhy{}\PYZhy{}\PYZhy{}\PYZhy{}\PYZhy{}\PYZhy{}\PYZhy{}\PYZhy{}\PYZhy{}\PYZhy{}\PYZhy{}\PYZhy{}\PYZhy{}\PYZhy{}\PYZhy{}\PYZhy{}\PYZhy{}\PYZhy{}}
\PY{c+c1}{\PYZsh{} Simulação e aplicação do PSIS}
\PY{c+c1}{\PYZsh{} \PYZhy{}\PYZhy{}\PYZhy{}\PYZhy{}\PYZhy{}\PYZhy{}\PYZhy{}\PYZhy{}\PYZhy{}\PYZhy{}\PYZhy{}\PYZhy{}\PYZhy{}\PYZhy{}\PYZhy{}\PYZhy{}\PYZhy{}\PYZhy{}\PYZhy{}\PYZhy{}\PYZhy{}\PYZhy{}\PYZhy{}\PYZhy{}\PYZhy{}\PYZhy{}\PYZhy{}\PYZhy{}\PYZhy{}\PYZhy{}\PYZhy{}\PYZhy{}\PYZhy{}\PYZhy{}\PYZhy{}\PYZhy{}\PYZhy{}\PYZhy{}\PYZhy{}\PYZhy{}\PYZhy{}\PYZhy{}\PYZhy{}\PYZhy{}\PYZhy{}\PYZhy{}\PYZhy{}\PYZhy{}\PYZhy{}\PYZhy{}\PYZhy{}\PYZhy{}\PYZhy{}\PYZhy{}\PYZhy{}\PYZhy{}\PYZhy{}\PYZhy{}\PYZhy{}\PYZhy{}\PYZhy{}\PYZhy{}\PYZhy{}\PYZhy{}\PYZhy{}\PYZhy{}\PYZhy{}\PYZhy{}\PYZhy{}\PYZhy{}\PYZhy{}\PYZhy{}\PYZhy{}}

\PY{n}{results} \PY{o}{=} \PY{p}{\PYZob{}}\PY{p}{\PYZcb{}}

\PY{k}{for} \PY{n}{scenario}\PY{p}{,} \PY{n}{spec} \PY{o+ow}{in} \PY{n}{scenarios}\PY{o}{.}\PY{n}{items}\PY{p}{(}\PY{p}{)}\PY{p}{:}

    \PY{n}{raw\PYZus{}log\PYZus{}weights} \PY{o}{=} \PY{n}{generate\PYZus{}log\PYZus{}weights}\PY{p}{(}\PY{n}{spec}\PY{p}{,} \PY{n}{n\PYZus{}samples}\PY{p}{,} \PY{n}{rng}\PY{p}{)}

    \PY{c+c1}{\PYZsh{} PSIS: log\PYZhy{}pesos suavizados (normalizados) e k\PYZhy{}hat}
    \PY{n}{smoothed\PYZus{}log\PYZus{}weights}\PY{p}{,} \PY{n}{k\PYZus{}hat} \PY{o}{=} \PY{n}{az}\PY{o}{.}\PY{n}{psislw}\PY{p}{(}\PY{n}{raw\PYZus{}log\PYZus{}weights}\PY{o}{.}\PY{n}{copy}\PY{p}{(}\PY{p}{)}\PY{p}{,} \PY{n}{reff}\PY{o}{=}\PY{n}{reff}\PY{p}{)}

    \PY{c+c1}{\PYZsh{} Pesos brutos normalizados e pesos suavizados}
    \PY{n}{raw\PYZus{}weights} \PY{o}{=} \PY{n}{np}\PY{o}{.}\PY{n}{exp}\PY{p}{(}\PY{n}{raw\PYZus{}log\PYZus{}weights} \PY{o}{\PYZhy{}} \PY{n}{logsumexp}\PY{p}{(}\PY{n}{raw\PYZus{}log\PYZus{}weights}\PY{p}{)}\PY{p}{)}
    \PY{n}{smoothed\PYZus{}weights} \PY{o}{=} \PY{n}{np}\PY{o}{.}\PY{n}{exp}\PY{p}{(}\PY{n}{smoothed\PYZus{}log\PYZus{}weights}\PY{p}{)}

    \PY{n}{results}\PY{p}{[}\PY{n}{scenario}\PY{p}{]} \PY{o}{=} \PY{p}{\PYZob{}}
        \PY{l+s+s2}{\PYZdq{}}\PY{l+s+s2}{k\PYZus{}true}\PY{l+s+s2}{\PYZdq{}}\PY{p}{:} \PY{n}{spec}\PY{p}{[}\PY{l+s+s2}{\PYZdq{}}\PY{l+s+s2}{k}\PY{l+s+s2}{\PYZdq{}}\PY{p}{]}\PY{p}{,}
        \PY{l+s+s2}{\PYZdq{}}\PY{l+s+s2}{k\PYZus{}hat}\PY{l+s+s2}{\PYZdq{}}\PY{p}{:} \PY{n+nb}{float}\PY{p}{(}\PY{n}{k\PYZus{}hat}\PY{p}{)}\PY{p}{,}
        \PY{l+s+s2}{\PYZdq{}}\PY{l+s+s2}{raw\PYZus{}sorted}\PY{l+s+s2}{\PYZdq{}}\PY{p}{:} \PY{n}{np}\PY{o}{.}\PY{n}{sort}\PY{p}{(}\PY{n}{raw\PYZus{}weights}\PY{p}{)}\PY{p}{,}
        \PY{l+s+s2}{\PYZdq{}}\PY{l+s+s2}{smoothed\PYZus{}sorted}\PY{l+s+s2}{\PYZdq{}}\PY{p}{:} \PY{n}{np}\PY{o}{.}\PY{n}{sort}\PY{p}{(}\PY{n}{smoothed\PYZus{}weights}\PY{p}{)}\PY{p}{,}
        \PY{l+s+s2}{\PYZdq{}}\PY{l+s+s2}{tail\PYZus{}length}\PY{l+s+s2}{\PYZdq{}}\PY{p}{:} \PY{n}{tail\PYZus{}length}\PY{p}{,}
    \PY{p}{\PYZcb{}}


\PY{c+c1}{\PYZsh{} \PYZhy{}\PYZhy{}\PYZhy{}\PYZhy{}\PYZhy{}\PYZhy{}\PYZhy{}\PYZhy{}\PYZhy{}\PYZhy{}\PYZhy{}\PYZhy{}\PYZhy{}\PYZhy{}\PYZhy{}\PYZhy{}\PYZhy{}\PYZhy{}\PYZhy{}\PYZhy{}\PYZhy{}\PYZhy{}\PYZhy{}\PYZhy{}\PYZhy{}\PYZhy{}\PYZhy{}\PYZhy{}\PYZhy{}\PYZhy{}\PYZhy{}\PYZhy{}\PYZhy{}\PYZhy{}\PYZhy{}\PYZhy{}\PYZhy{}\PYZhy{}\PYZhy{}\PYZhy{}\PYZhy{}\PYZhy{}\PYZhy{}\PYZhy{}\PYZhy{}\PYZhy{}\PYZhy{}\PYZhy{}\PYZhy{}\PYZhy{}\PYZhy{}\PYZhy{}\PYZhy{}\PYZhy{}\PYZhy{}\PYZhy{}\PYZhy{}\PYZhy{}\PYZhy{}\PYZhy{}\PYZhy{}\PYZhy{}\PYZhy{}\PYZhy{}\PYZhy{}\PYZhy{}\PYZhy{}\PYZhy{}\PYZhy{}\PYZhy{}\PYZhy{}\PYZhy{}\PYZhy{}}
\PY{c+c1}{\PYZsh{} Visualização}
\PY{c+c1}{\PYZsh{} \PYZhy{}\PYZhy{}\PYZhy{}\PYZhy{}\PYZhy{}\PYZhy{}\PYZhy{}\PYZhy{}\PYZhy{}\PYZhy{}\PYZhy{}\PYZhy{}\PYZhy{}\PYZhy{}\PYZhy{}\PYZhy{}\PYZhy{}\PYZhy{}\PYZhy{}\PYZhy{}\PYZhy{}\PYZhy{}\PYZhy{}\PYZhy{}\PYZhy{}\PYZhy{}\PYZhy{}\PYZhy{}\PYZhy{}\PYZhy{}\PYZhy{}\PYZhy{}\PYZhy{}\PYZhy{}\PYZhy{}\PYZhy{}\PYZhy{}\PYZhy{}\PYZhy{}\PYZhy{}\PYZhy{}\PYZhy{}\PYZhy{}\PYZhy{}\PYZhy{}\PYZhy{}\PYZhy{}\PYZhy{}\PYZhy{}\PYZhy{}\PYZhy{}\PYZhy{}\PYZhy{}\PYZhy{}\PYZhy{}\PYZhy{}\PYZhy{}\PYZhy{}\PYZhy{}\PYZhy{}\PYZhy{}\PYZhy{}\PYZhy{}\PYZhy{}\PYZhy{}\PYZhy{}\PYZhy{}\PYZhy{}\PYZhy{}\PYZhy{}\PYZhy{}\PYZhy{}\PYZhy{}}

\PY{n}{fig}\PY{p}{,} \PY{n}{axes} \PY{o}{=} \PY{n}{plt}\PY{o}{.}\PY{n}{subplots}\PY{p}{(}
    \PY{n}{nrows}\PY{o}{=}\PY{n+nb}{len}\PY{p}{(}\PY{n}{results}\PY{p}{)}\PY{p}{,} \PY{n}{ncols}\PY{o}{=}\PY{l+m+mi}{3}\PY{p}{,}
    \PY{n}{figsize}\PY{o}{=}\PY{p}{(}\PY{l+m+mi}{30}\PY{p}{,} \PY{l+m+mi}{24}\PY{p}{)}
\PY{p}{)}

\PY{n}{fig}\PY{o}{.}\PY{n}{suptitle}\PY{p}{(}
    \PY{l+s+s2}{\PYZdq{}}\PY{l+s+s2}{Suavização de Pareto dos pesos de importância}\PY{l+s+s2}{\PYZdq{}}\PY{p}{,}
    \PY{n}{fontsize}\PY{o}{=}\PY{l+m+mi}{24}\PY{p}{,}
    \PY{n}{y}\PY{o}{=}\PY{l+m+mf}{0.995}
\PY{p}{)}

\PY{n}{ranks} \PY{o}{=} \PY{n}{np}\PY{o}{.}\PY{n}{arange}\PY{p}{(}\PY{l+m+mi}{1}\PY{p}{,} \PY{n}{n\PYZus{}samples} \PY{o}{+} \PY{l+m+mi}{1}\PY{p}{)}   \PY{c+c1}{\PYZsh{} posição a partir do menor peso}

\PY{k}{for} \PY{n}{row}\PY{p}{,} \PY{p}{(}\PY{p}{(}\PY{n}{scenario}\PY{p}{,} \PY{n}{r}\PY{p}{)}\PY{p}{,} \PY{n}{color}\PY{p}{)} \PY{o+ow}{in} \PY{n+nb}{enumerate}\PY{p}{(}\PY{n+nb}{zip}\PY{p}{(}\PY{n}{results}\PY{o}{.}\PY{n}{items}\PY{p}{(}\PY{p}{)}\PY{p}{,} \PY{n}{colors}\PY{p}{)}\PY{p}{)}\PY{p}{:}

    \PY{n}{ax\PYZus{}w}\PY{p}{,} \PY{n}{ax\PYZus{}tail}\PY{p}{,} \PY{n}{ax\PYZus{}ratio} \PY{o}{=} \PY{n}{axes}\PY{p}{[}\PY{n}{row}\PY{p}{]}
    \PY{n}{M} \PY{o}{=} \PY{n}{r}\PY{p}{[}\PY{l+s+s2}{\PYZdq{}}\PY{l+s+s2}{tail\PYZus{}length}\PY{l+s+s2}{\PYZdq{}}\PY{p}{]}

    \PY{c+c1}{\PYZsh{} \PYZhy{}\PYZhy{}\PYZhy{}\PYZhy{} Painel esquerdo: pesos ordenados do menor para o maior (log\PYZhy{}log) \PYZhy{}}
    \PY{n}{ax\PYZus{}w}\PY{o}{.}\PY{n}{scatter}\PY{p}{(}
        \PY{n}{ranks}\PY{p}{,} \PY{n}{r}\PY{p}{[}\PY{l+s+s2}{\PYZdq{}}\PY{l+s+s2}{raw\PYZus{}sorted}\PY{l+s+s2}{\PYZdq{}}\PY{p}{]}\PY{p}{,}
        \PY{n}{s}\PY{o}{=}\PY{l+m+mi}{18}\PY{p}{,} \PY{n}{color}\PY{o}{=}\PY{n}{color}\PY{p}{,} \PY{n}{alpha}\PY{o}{=}\PY{l+m+mf}{0.5}\PY{p}{,}
        \PY{n}{label}\PY{o}{=}\PY{l+s+s2}{\PYZdq{}}\PY{l+s+s2}{Pesos brutos}\PY{l+s+s2}{\PYZdq{}}
    \PY{p}{)}
    \PY{n}{ax\PYZus{}w}\PY{o}{.}\PY{n}{plot}\PY{p}{(}
        \PY{n}{ranks}\PY{p}{,} \PY{n}{r}\PY{p}{[}\PY{l+s+s2}{\PYZdq{}}\PY{l+s+s2}{smoothed\PYZus{}sorted}\PY{l+s+s2}{\PYZdq{}}\PY{p}{]}\PY{p}{,}
        \PY{n}{color}\PY{o}{=}\PY{n}{color}\PY{p}{,} \PY{n}{linewidth}\PY{o}{=}\PY{l+m+mi}{3}\PY{p}{,}
        \PY{n}{label}\PY{o}{=}\PY{l+s+s2}{\PYZdq{}}\PY{l+s+s2}{Pesos suavizados (PSIS)}\PY{l+s+s2}{\PYZdq{}}
    \PY{p}{)}
    \PY{n}{ax\PYZus{}w}\PY{o}{.}\PY{n}{axvline}\PY{p}{(}
        \PY{n}{n\PYZus{}samples} \PY{o}{\PYZhy{}} \PY{n}{M}\PY{p}{,}
        \PY{n}{color}\PY{o}{=}\PY{l+s+s2}{\PYZdq{}}\PY{l+s+s2}{black}\PY{l+s+s2}{\PYZdq{}}\PY{p}{,} \PY{n}{linestyle}\PY{o}{=}\PY{l+s+s2}{\PYZdq{}}\PY{l+s+s2}{\PYZhy{}\PYZhy{}}\PY{l+s+s2}{\PYZdq{}}\PY{p}{,} \PY{n}{linewidth}\PY{o}{=}\PY{l+m+mf}{1.8}\PY{p}{,}
        \PY{n}{label}\PY{o}{=}\PY{l+s+sa}{f}\PY{l+s+s2}{\PYZdq{}}\PY{l+s+s2}{Início da cauda suavizada (M=}\PY{l+s+si}{\PYZob{}}\PY{n}{M}\PY{l+s+si}{\PYZcb{}}\PY{l+s+s2}{)}\PY{l+s+s2}{\PYZdq{}}
    \PY{p}{)}

    \PY{n}{ax\PYZus{}w}\PY{o}{.}\PY{n}{set\PYZus{}xscale}\PY{p}{(}\PY{l+s+s2}{\PYZdq{}}\PY{l+s+s2}{log}\PY{l+s+s2}{\PYZdq{}}\PY{p}{)}
    \PY{k}{if} \PY{n}{log\PYZus{}y}\PY{p}{:}
        \PY{n}{ax\PYZus{}w}\PY{o}{.}\PY{n}{set\PYZus{}yscale}\PY{p}{(}\PY{l+s+s2}{\PYZdq{}}\PY{l+s+s2}{log}\PY{l+s+s2}{\PYZdq{}}\PY{p}{)}

    \PY{n}{ax\PYZus{}w}\PY{o}{.}\PY{n}{set\PYZus{}xlabel}\PY{p}{(}\PY{l+s+s2}{\PYZdq{}}\PY{l+s+s2}{Posição a partir do menor peso}\PY{l+s+s2}{\PYZdq{}}\PY{p}{,} \PY{n}{fontsize}\PY{o}{=}\PY{l+m+mi}{15}\PY{p}{)}
    \PY{n}{ax\PYZus{}w}\PY{o}{.}\PY{n}{set\PYZus{}ylabel}\PY{p}{(}\PY{l+s+s2}{\PYZdq{}}\PY{l+s+s2}{Peso normalizado}\PY{l+s+s2}{\PYZdq{}}\PY{p}{,} \PY{n}{fontsize}\PY{o}{=}\PY{l+m+mi}{15}\PY{p}{)}
    \PY{n}{ax\PYZus{}w}\PY{o}{.}\PY{n}{set\PYZus{}title}\PY{p}{(}
        \PY{l+s+s2}{\PYZdq{}}\PY{l+s+s2}{Pesos normalizados e ponto de corte: }\PY{l+s+s2}{\PYZdq{}}
        \PY{l+s+sa}{rf}\PY{l+s+s2}{\PYZdq{}}\PY{l+s+s2}{\PYZdl{}}\PY{l+s+s2}{\PYZbs{}}\PY{l+s+s2}{widehat}\PY{l+s+se}{\PYZob{}\PYZob{}}\PY{l+s+s2}{k}\PY{l+s+se}{\PYZcb{}\PYZcb{}}\PY{l+s+s2}{=}\PY{l+s+si}{\PYZob{}}\PY{n}{r}\PY{p}{[}\PY{l+s+s1}{\PYZsq{}}\PY{l+s+s1}{k\PYZus{}hat}\PY{l+s+s1}{\PYZsq{}}\PY{p}{]}\PY{l+s+si}{:}\PY{l+s+s2}{.2f}\PY{l+s+si}{\PYZcb{}}\PY{l+s+s2}{\PYZdl{}}\PY{l+s+s2}{\PYZdq{}}\PY{p}{,}
        \PY{n}{fontsize}\PY{o}{=}\PY{l+m+mi}{17}
    \PY{p}{)}
    \PY{n}{ax\PYZus{}w}\PY{o}{.}\PY{n}{tick\PYZus{}params}\PY{p}{(}\PY{n}{axis}\PY{o}{=}\PY{l+s+s2}{\PYZdq{}}\PY{l+s+s2}{both}\PY{l+s+s2}{\PYZdq{}}\PY{p}{,} \PY{n}{labelsize}\PY{o}{=}\PY{l+m+mi}{13}\PY{p}{)}
    \PY{n}{ax\PYZus{}w}\PY{o}{.}\PY{n}{grid}\PY{p}{(}\PY{n}{alpha}\PY{o}{=}\PY{l+m+mf}{0.25}\PY{p}{,} \PY{n}{linestyle}\PY{o}{=}\PY{l+s+s2}{\PYZdq{}}\PY{l+s+s2}{\PYZhy{}\PYZhy{}}\PY{l+s+s2}{\PYZdq{}}\PY{p}{)}
    \PY{n}{ax\PYZus{}w}\PY{o}{.}\PY{n}{legend}\PY{p}{(}\PY{n}{fontsize}\PY{o}{=}\PY{l+m+mi}{12}\PY{p}{,} \PY{n}{loc}\PY{o}{=}\PY{l+s+s2}{\PYZdq{}}\PY{l+s+s2}{upper left}\PY{l+s+s2}{\PYZdq{}}\PY{p}{)}

    \PY{c+c1}{\PYZsh{} \PYZhy{}\PYZhy{}\PYZhy{}\PYZhy{} Painel central: mesmos pesos, restrito à cauda (eixo x linear) \PYZhy{}\PYZhy{}}
    \PY{n}{tail\PYZus{}ranks} \PY{o}{=} \PY{n}{ranks}\PY{p}{[}\PY{o}{\PYZhy{}}\PY{n}{M}\PY{p}{:}\PY{p}{]}

    \PY{n}{ax\PYZus{}tail}\PY{o}{.}\PY{n}{scatter}\PY{p}{(}
        \PY{n}{tail\PYZus{}ranks}\PY{p}{,} \PY{n}{r}\PY{p}{[}\PY{l+s+s2}{\PYZdq{}}\PY{l+s+s2}{raw\PYZus{}sorted}\PY{l+s+s2}{\PYZdq{}}\PY{p}{]}\PY{p}{[}\PY{o}{\PYZhy{}}\PY{n}{M}\PY{p}{:}\PY{p}{]}\PY{p}{,}
        \PY{n}{s}\PY{o}{=}\PY{l+m+mi}{24}\PY{p}{,} \PY{n}{color}\PY{o}{=}\PY{n}{color}\PY{p}{,} \PY{n}{alpha}\PY{o}{=}\PY{l+m+mf}{0.5}\PY{p}{,}
        \PY{n}{label}\PY{o}{=}\PY{l+s+s2}{\PYZdq{}}\PY{l+s+s2}{Pesos brutos}\PY{l+s+s2}{\PYZdq{}}
    \PY{p}{)}
    \PY{n}{ax\PYZus{}tail}\PY{o}{.}\PY{n}{plot}\PY{p}{(}
        \PY{n}{tail\PYZus{}ranks}\PY{p}{,} \PY{n}{r}\PY{p}{[}\PY{l+s+s2}{\PYZdq{}}\PY{l+s+s2}{smoothed\PYZus{}sorted}\PY{l+s+s2}{\PYZdq{}}\PY{p}{]}\PY{p}{[}\PY{o}{\PYZhy{}}\PY{n}{M}\PY{p}{:}\PY{p}{]}\PY{p}{,}
        \PY{n}{color}\PY{o}{=}\PY{n}{color}\PY{p}{,} \PY{n}{linewidth}\PY{o}{=}\PY{l+m+mi}{3}\PY{p}{,}
        \PY{n}{label}\PY{o}{=}\PY{l+s+s2}{\PYZdq{}}\PY{l+s+s2}{Pesos suavizados (PSIS)}\PY{l+s+s2}{\PYZdq{}}
    \PY{p}{)}

    \PY{k}{if} \PY{n}{log\PYZus{}y}\PY{p}{:}
        \PY{n}{ax\PYZus{}tail}\PY{o}{.}\PY{n}{set\PYZus{}yscale}\PY{p}{(}\PY{l+s+s2}{\PYZdq{}}\PY{l+s+s2}{log}\PY{l+s+s2}{\PYZdq{}}\PY{p}{)}

    \PY{n}{ax\PYZus{}tail}\PY{o}{.}\PY{n}{set\PYZus{}xlabel}\PY{p}{(}\PY{l+s+s2}{\PYZdq{}}\PY{l+s+s2}{Posição a partir do menor peso}\PY{l+s+s2}{\PYZdq{}}\PY{p}{,} \PY{n}{fontsize}\PY{o}{=}\PY{l+m+mi}{15}\PY{p}{)}
    \PY{n}{ax\PYZus{}tail}\PY{o}{.}\PY{n}{set\PYZus{}ylabel}\PY{p}{(}\PY{l+s+s2}{\PYZdq{}}\PY{l+s+s2}{Peso normalizado}\PY{l+s+s2}{\PYZdq{}}\PY{p}{,} \PY{n}{fontsize}\PY{o}{=}\PY{l+m+mi}{15}\PY{p}{)}
    \PY{n}{ax\PYZus{}tail}\PY{o}{.}\PY{n}{set\PYZus{}title}\PY{p}{(}
        \PY{l+s+s2}{\PYZdq{}}\PY{l+s+s2}{Pesos na cauda suavizada: }\PY{l+s+s2}{\PYZdq{}}
        \PY{l+s+sa}{rf}\PY{l+s+s2}{\PYZdq{}}\PY{l+s+s2}{\PYZdl{}}\PY{l+s+s2}{\PYZbs{}}\PY{l+s+s2}{widehat}\PY{l+s+se}{\PYZob{}\PYZob{}}\PY{l+s+s2}{k}\PY{l+s+se}{\PYZcb{}\PYZcb{}}\PY{l+s+s2}{=}\PY{l+s+si}{\PYZob{}}\PY{n}{r}\PY{p}{[}\PY{l+s+s1}{\PYZsq{}}\PY{l+s+s1}{k\PYZus{}hat}\PY{l+s+s1}{\PYZsq{}}\PY{p}{]}\PY{l+s+si}{:}\PY{l+s+s2}{.2f}\PY{l+s+si}{\PYZcb{}}\PY{l+s+s2}{\PYZdl{}}\PY{l+s+s2}{\PYZdq{}}\PY{p}{,}
        \PY{n}{fontsize}\PY{o}{=}\PY{l+m+mi}{17}
    \PY{p}{)}
    \PY{n}{ax\PYZus{}tail}\PY{o}{.}\PY{n}{tick\PYZus{}params}\PY{p}{(}\PY{n}{axis}\PY{o}{=}\PY{l+s+s2}{\PYZdq{}}\PY{l+s+s2}{both}\PY{l+s+s2}{\PYZdq{}}\PY{p}{,} \PY{n}{labelsize}\PY{o}{=}\PY{l+m+mi}{13}\PY{p}{)}
    \PY{n}{ax\PYZus{}tail}\PY{o}{.}\PY{n}{grid}\PY{p}{(}\PY{n}{alpha}\PY{o}{=}\PY{l+m+mf}{0.25}\PY{p}{,} \PY{n}{linestyle}\PY{o}{=}\PY{l+s+s2}{\PYZdq{}}\PY{l+s+s2}{\PYZhy{}\PYZhy{}}\PY{l+s+s2}{\PYZdq{}}\PY{p}{)}
    \PY{n}{ax\PYZus{}tail}\PY{o}{.}\PY{n}{legend}\PY{p}{(}\PY{n}{fontsize}\PY{o}{=}\PY{l+m+mi}{12}\PY{p}{,} \PY{n}{loc}\PY{o}{=}\PY{l+s+s2}{\PYZdq{}}\PY{l+s+s2}{upper left}\PY{l+s+s2}{\PYZdq{}}\PY{p}{)}

    \PY{c+c1}{\PYZsh{} \PYZhy{}\PYZhy{}\PYZhy{}\PYZhy{} Painel direito: razão suavizado / bruto na cauda \PYZhy{}\PYZhy{}\PYZhy{}\PYZhy{}\PYZhy{}\PYZhy{}\PYZhy{}\PYZhy{}\PYZhy{}\PYZhy{}\PYZhy{}\PYZhy{}\PYZhy{}\PYZhy{}\PYZhy{}}
    \PY{n}{ax\PYZus{}ratio}\PY{o}{.}\PY{n}{axhline}\PY{p}{(}\PY{l+m+mi}{1}\PY{p}{,} \PY{n}{color}\PY{o}{=}\PY{l+s+s2}{\PYZdq{}}\PY{l+s+s2}{lightgray}\PY{l+s+s2}{\PYZdq{}}\PY{p}{,} \PY{n}{linestyle}\PY{o}{=}\PY{l+s+s2}{\PYZdq{}}\PY{l+s+s2}{\PYZhy{}\PYZhy{}}\PY{l+s+s2}{\PYZdq{}}\PY{p}{)}
    \PY{n}{ax\PYZus{}ratio}\PY{o}{.}\PY{n}{plot}\PY{p}{(}
        \PY{n}{tail\PYZus{}ranks}\PY{p}{,} \PY{n}{r}\PY{p}{[}\PY{l+s+s2}{\PYZdq{}}\PY{l+s+s2}{smoothed\PYZus{}sorted}\PY{l+s+s2}{\PYZdq{}}\PY{p}{]}\PY{p}{[}\PY{o}{\PYZhy{}}\PY{n}{M}\PY{p}{:}\PY{p}{]} \PY{o}{/} \PY{n}{r}\PY{p}{[}\PY{l+s+s2}{\PYZdq{}}\PY{l+s+s2}{raw\PYZus{}sorted}\PY{l+s+s2}{\PYZdq{}}\PY{p}{]}\PY{p}{[}\PY{o}{\PYZhy{}}\PY{n}{M}\PY{p}{:}\PY{p}{]}\PY{p}{,}
        \PY{n}{color}\PY{o}{=}\PY{n}{color}\PY{p}{,} \PY{n}{linewidth}\PY{o}{=}\PY{l+m+mf}{2.5}
    \PY{p}{)}

    \PY{n}{ax\PYZus{}ratio}\PY{o}{.}\PY{n}{set\PYZus{}xlabel}\PY{p}{(}\PY{l+s+s2}{\PYZdq{}}\PY{l+s+s2}{Posição a partir do menor peso}\PY{l+s+s2}{\PYZdq{}}\PY{p}{,} \PY{n}{fontsize}\PY{o}{=}\PY{l+m+mi}{15}\PY{p}{)}
    \PY{n}{ax\PYZus{}ratio}\PY{o}{.}\PY{n}{set\PYZus{}ylabel}\PY{p}{(}\PY{l+s+s2}{\PYZdq{}}\PY{l+s+s2}{Suavizado / bruto}\PY{l+s+s2}{\PYZdq{}}\PY{p}{,} \PY{n}{fontsize}\PY{o}{=}\PY{l+m+mi}{15}\PY{p}{)}
    \PY{n}{ax\PYZus{}ratio}\PY{o}{.}\PY{n}{set\PYZus{}title}\PY{p}{(}
        \PY{l+s+s2}{\PYZdq{}}\PY{l+s+s2}{Efeito da suavização na cauda: }\PY{l+s+s2}{\PYZdq{}}
        \PY{l+s+sa}{rf}\PY{l+s+s2}{\PYZdq{}}\PY{l+s+s2}{\PYZdl{}}\PY{l+s+s2}{\PYZbs{}}\PY{l+s+s2}{widehat}\PY{l+s+se}{\PYZob{}\PYZob{}}\PY{l+s+s2}{k}\PY{l+s+se}{\PYZcb{}\PYZcb{}}\PY{l+s+s2}{=}\PY{l+s+si}{\PYZob{}}\PY{n}{r}\PY{p}{[}\PY{l+s+s1}{\PYZsq{}}\PY{l+s+s1}{k\PYZus{}hat}\PY{l+s+s1}{\PYZsq{}}\PY{p}{]}\PY{l+s+si}{:}\PY{l+s+s2}{.2f}\PY{l+s+si}{\PYZcb{}}\PY{l+s+s2}{\PYZdl{}}\PY{l+s+s2}{\PYZdq{}}\PY{p}{,}
        \PY{n}{fontsize}\PY{o}{=}\PY{l+m+mi}{17}
    \PY{p}{)}
    \PY{n}{ax\PYZus{}ratio}\PY{o}{.}\PY{n}{tick\PYZus{}params}\PY{p}{(}\PY{n}{axis}\PY{o}{=}\PY{l+s+s2}{\PYZdq{}}\PY{l+s+s2}{both}\PY{l+s+s2}{\PYZdq{}}\PY{p}{,} \PY{n}{labelsize}\PY{o}{=}\PY{l+m+mi}{13}\PY{p}{)}
    \PY{n}{ax\PYZus{}ratio}\PY{o}{.}\PY{n}{grid}\PY{p}{(}\PY{n}{alpha}\PY{o}{=}\PY{l+m+mf}{0.25}\PY{p}{,} \PY{n}{linestyle}\PY{o}{=}\PY{l+s+s2}{\PYZdq{}}\PY{l+s+s2}{\PYZhy{}\PYZhy{}}\PY{l+s+s2}{\PYZdq{}}\PY{p}{)}

\PY{n}{plt}\PY{o}{.}\PY{n}{tight\PYZus{}layout}\PY{p}{(}\PY{n}{rect}\PY{o}{=}\PY{p}{[}\PY{l+m+mi}{0}\PY{p}{,} \PY{l+m+mi}{0}\PY{p}{,} \PY{l+m+mi}{1}\PY{p}{,} \PY{l+m+mf}{0.98}\PY{p}{]}\PY{p}{)}
\PY{n}{plt}\PY{o}{.}\PY{n}{show}\PY{p}{(}\PY{p}{)}
\end{Verbatim}
\end{tcolorbox}

    \begin{Verbatim}[commandchars=\\\{\}]
\#\#\#\#\#\#\#\#\#\#\#\#\#\#\#\#\#\#\#\#\#\#\#\#\#\#\#\#\#\#\#\#\#\#\#\#\#\#\#\#\#\#\#\#\#\#\#\#\#\#\#\#\#\#\#\#\#\#\#\#\#\#\#\#\#\#\#\#\#\#\#\#\#\#\#\#\#\#\#\#
\#\#\#\#\#\#\#\#\#\#\#\#\#\#\#\#\#\#\#\#
CONFIGURAÇÃO DA SIMULAÇÃO: PSIS SOBRE PESOS DE IMPORTÂNCIA SINTÉTICOS
\#\#\#\#\#\#\#\#\#\#\#\#\#\#\#\#\#\#\#\#\#\#\#\#\#\#\#\#\#\#\#\#\#\#\#\#\#\#\#\#\#\#\#\#\#\#\#\#\#\#\#\#\#\#\#\#\#\#\#\#\#\#\#\#\#\#\#\#\#\#\#\#\#\#\#\#\#\#\#\#
\#\#\#\#\#\#\#\#\#\#\#\#\#\#\#\#\#\#\#\#
Amostras:         K = 5000  (reff = 1.0)
Cauda do PSIS:    M = 213 maiores pesos (4.3\% da amostra)
Suavização:       Pareto generalizada ajustada à cauda (arviz.psislw)
--------------------------------------------------------------------------------
--------------------
Geração dos pesos brutos (forma funcional única):
  log w = k * E, com E \textasciitilde{} Exponencial(1)
  Escala maior (decaimento mais lento) -> cauda mais pesada -> k maior.
  Escala menor (decaimento mais rápido) -> cauda mais leve -> k menor.
\#\#\#\#\#\#\#\#\#\#\#\#\#\#\#\#\#\#\#\#\#\#\#\#\#\#\#\#\#\#\#\#\#\#\#\#\#\#\#\#\#\#\#\#\#\#\#\#\#\#\#\#\#\#\#\#\#\#\#\#\#\#\#\#\#\#\#\#\#\#\#\#\#\#\#\#\#\#\#\#
\#\#\#\#\#\#\#\#\#\#\#\#\#\#\#\#\#\#\#\#
Leitura: k-hat < 0,5 (leve); 0,5-0,7 (limiar); > 0,7 (pouco confiável).

    \end{Verbatim}

    \begin{center}
    \adjustimage{max size={0.9\linewidth}{0.9\paperheight}}{notebook_corrigido_files/notebook_corrigido_5_1.png}
    \end{center}
    { \hspace*{\fill} \\}
    
    \subsubsection{8.2 Sensibilidade local (Pérez, Martín \& Rufo,
2006)}\label{sensibilidade-local-puxe9rez-martuxedn-rufo-2006}

As derivadas parciais de uma quantidade a posteriori em relação ao
parâmetro perturbado podem ser estimadas por Monte Carlo usando as
próprias amostras da posteriori base, sem nova amostragem.

\paragraph{8.2.1 Identidade geral e estimador
MCMC}\label{identidade-geral-e-estimador-mcmc}

Seja \(\mathcal H(\xi)=\mathbb E_{\xi}[h(\mathbf B)]\), em que \(\xi\)
representa os parâmetros da priori e do modelo das fontes, e \(\xi_j\) é
o componente perturbado. Considerando \(h\) fixa, sem dependência
explícita de \(\xi\), tem-se:

\[\left[
\frac{\partial \mathcal H}{\partial \xi_j}
\right]_{\xi=\xi^{0}}
=
\mathbb E_{0}
\left[
\big(h(\mathbf B)-\mathcal H^{0}\big)
\eta_{\xi_j}(\mathbf B)
\right]
=
\operatorname{Cov}_{0}
\big(h,\eta_{\xi_j}\big).\]

O índice \(0\) indica que as esperanças e covariâncias são calculadas
sob a posteriori base.

O escore da posteriori é a derivada parcial de sua log-densidade em
relação ao parâmetro perturbado:

\[\eta_{\xi_j}(\mathbf B)
=
\left[
\frac{\partial
\log p_{\mathbf B\mid\mathcal X}
(\mathbf B\mid\mathcal X;\xi)}
{\partial \xi_j}
\right]_{\xi=\xi^{0}}.\]

Para perturbações da priori ou da verossimilhança, podem-se usar os
escores abaixo. Eles diferem do escore da posteriori normalizada apenas
por uma constante em \(\mathbf B\), que não altera a covariância.

{\def\LTcaptype{none} % do not increment counter
\begin{longtable}[]{@{}
  >{\raggedright\arraybackslash}p{(\linewidth - 2\tabcolsep) * \real{0.5000}}
  >{\raggedright\arraybackslash}p{(\linewidth - 2\tabcolsep) * \real{0.5000}}@{}}
\toprule\noalign{}
\begin{minipage}[b]{\linewidth}\raggedright
Perturbação
\end{minipage} & \begin{minipage}[b]{\linewidth}\raggedright
Escore equivalente na covariância
\end{minipage} \\
\midrule\noalign{}
\endhead
\bottomrule\noalign{}
\endlastfoot
Priori &
\(\eta_{\xi_p}(\mathbf B)=\left[\frac{\partial \log p_{\mathbf b}(\mathbf b;\xi_p)}{\partial \xi_p}\right]_{\xi_p=\xi_p^{0}}\),
com \(\mathbf b=\mathrm{vec}(\mathbf B)\) \\
Modelo das fontes &
\(\eta_{\xi_s}(\mathbf B)=\left[\frac{\partial \log \mathcal L_{\mathcal X}(\mathbf B;\xi_s)}{\partial \xi_s}\right]_{\xi_s=\xi_s^{0}}\) \\
\end{longtable}
}

Com \(\hat{\mathcal H}^{0}=K^{-1}\sum_{k=1}^{K}h(\mathbf B_k)\), o
estimador MCMC correspondente é:

\[\widehat{
\left[
\frac{\partial \mathcal H}{\partial \xi_j}
\right]_{\xi=\xi^{0}}
}
=
\frac{1}{K}
\sum_{k=1}^{K}
\big(h(\mathbf B_k)-\hat{\mathcal H}^{0}\big)
\eta_{\xi_j}(\mathbf B_k).\]

\paragraph{8.2.2 Aplicação à média
posterior}\label{aplicauxe7uxe3o-uxe0-muxe9dia-posterior}

Neste ensaio, a análise local considera a média posterior de cada
elemento \(b_{ij}\) da matriz de separação:

\[\bar b_{ij}(\xi)=\mathbb E_{\xi}[b_{ij}].\]

Como \(h(\mathbf B)=b_{ij}\) não depende explicitamente de \(\xi\), sua
sensibilidade é:

\[\left[
\frac{\partial \bar b_{ij}}{\partial \xi_j}
\right]_{\xi=\xi^{0}}
=
\operatorname{Cov}_{0}(b_{ij},\eta_{\xi_j}).\]

Para evitar confusão entre o índice da matriz e o índice do parâmetro
perturbado, podemos escrever este último como \(a\):

\[\left[
\frac{\partial \bar b_{ij}}{\partial \xi_a}
\right]_{\xi=\xi^{0}}
=
\operatorname{Cov}_{0}(b_{ij},\eta_{\xi_a}).\]

Com \(\hat{\bar b}_{ij}^{0}=K^{-1}\sum_{k=1}^{K}[\mathbf B_k]_{ij}\), a
estimativa é:

\[\widehat{
\left[
\frac{\partial \bar b_{ij}}{\partial \xi_a}
\right]_{\xi=\xi^{0}}
}
=
\frac{1}{K}
\sum_{k=1}^{K}
\big([\mathbf B_k]_{ij}-\hat{\bar b}_{ij}^{0}\big)
\eta_{\xi_a}(\mathbf B_k).\]

Essa derivada quantifica a mudança local da média posterior em resposta
à variação de um parâmetro da priori ou do modelo das fontes, avaliada
no ponto base.

\subsubsection{8.3 Escores usados}\label{escores-usados}

As expressões abaixo são derivadas das log-densidades em relação aos
parâmetros diretamente, mantendo os demais fixos. Para a sensibilidade
local, todos os parâmetros são avaliados nos valores base \(\xi_p^{0}\)
e \(\xi_s^{0}\).

Considera-se uma Student-t de localização \(\mu_s\) e escala
\(\sigma_s\) para as fontes, e uma t multivariada de localização
\(\boldsymbol\theta_p\) e matriz de forma \(\sigma_p^2\mathbf I_d\) para
a priori.

Notação:

\[z=\frac{[\mathbf B\mathbf x(n)]_\ell-\mu_s}{\sigma_s},
\qquad
\mathbf b=\mathrm{vec}(\mathbf B),
\qquad
\delta^2=
\frac{(\mathbf b-\boldsymbol\theta_p)^\top
(\mathbf b-\boldsymbol\theta_p)}{\sigma_p^2},
\qquad d=L^2=4.\]

Os escores das fontes apresentados na tabela correspondem a uma
observação e uma fonte. Para obter o escore da verossimilhança completa,
somam-se as contribuições em \(n\) e \(\ell\).

{\def\LTcaptype{none} % do not increment counter
\begin{longtable}[]{@{}
  >{\raggedright\arraybackslash}p{(\linewidth - 2\tabcolsep) * \real{0.5000}}
  >{\raggedright\arraybackslash}p{(\linewidth - 2\tabcolsep) * \real{0.5000}}@{}}
\toprule\noalign{}
\begin{minipage}[b]{\linewidth}\raggedright
Parâmetro
\end{minipage} & \begin{minipage}[b]{\linewidth}\raggedright
Escore
\end{minipage} \\
\midrule\noalign{}
\endhead
\bottomrule\noalign{}
\endlastfoot
\(\mu_s\) & \(\dfrac{(\nu_s+1)z}{\sigma_s(\nu_s+z^2)}\) \\
\(\sigma_s\) &
\(\dfrac{1}{\sigma_s}\left[-1+\dfrac{(\nu_s+1)z^2}{\nu_s+z^2}\right]\) \\
\(\nu_s\) &
\(\dfrac12\left[\psi\!\left(\dfrac{\nu_s+1}{2}\right)-\psi\!\left(\dfrac{\nu_s}{2}\right)-\dfrac1{\nu_s}-\log\!\left(1+\dfrac{z^2}{\nu_s}\right)+\dfrac{(\nu_s+1)z^2}{\nu_s(\nu_s+z^2)}\right]\) \\
\(\alpha\), com
\(\partial\boldsymbol\theta_p/\partial\alpha=\mathbf i-\underline{\mathbf b}\)
&
\(\dfrac{(\nu_p+d)(\mathbf b-\boldsymbol\theta_p)^\top(\mathbf i-\underline{\mathbf b})}{\nu_p\sigma_p^2+(\mathbf b-\boldsymbol\theta_p)^\top(\mathbf b-\boldsymbol\theta_p)}\) \\
\(\sigma_p\) &
\(\dfrac{1}{\sigma_p}\left[-d+\dfrac{(\nu_p+d)\delta^2}{\nu_p+\delta^2}\right]\) \\
\(\nu_p\) &
\(\dfrac12\left[\psi\!\left(\dfrac{\nu_p+d}{2}\right)-\psi\!\left(\dfrac{\nu_p}{2}\right)-\dfrac d{\nu_p}-\log\!\left(1+\dfrac{\delta^2}{\nu_p}\right)+\dfrac{(\nu_p+d)\delta^2}{\nu_p(\nu_p+\delta^2)}\right]\) \\
\end{longtable}
}

A função \(\psi\), denominada digama, é definida por:

\[\psi(t) = \frac{d}{dt}\log\Gamma(t) = \frac{\Gamma'(t)}{\Gamma(t)}.\]

\subsection{9. Saídas esperadas (por fonte × perturbação ×
realização)}\label{sauxeddas-esperadas-por-fonte-perturbauxe7uxe3o-realizauxe7uxe3o}

\begin{itemize}
\tightlist
\item
  Curvas dos 4 momentos de \(b_{ij}\) em função do parâmetro perturbado,
  considerando tanto hiperparâmetros da priori
  (\(\boldsymbol\theta_p^{*}\), \(\sigma_p^{*}\) e \(\nu_p^{*}\)) quanto
  do modelo das fontes (\(\mu_s^{*}\), \(\sigma_s^{*}\) e
  \(\nu_s^{*}\)).
\item
  Derivadas locais da média de \(b_{ij}\), no ponto base.
\item
  Curva de \(\hat k\) → região de validade.
\item
  Agregação nas realizações que convergiram, comparando rs1 → rs10.
\end{itemize}

\subsection{10. Posicionamento metodológico (1
slide)}\label{posicionamento-metodoluxf3gico-1-slide}

{\def\LTcaptype{none} % do not increment counter
\begin{longtable}[]{@{}
  >{\raggedright\arraybackslash}p{(\linewidth - 4\tabcolsep) * \real{0.3333}}
  >{\raggedright\arraybackslash}p{(\linewidth - 4\tabcolsep) * \real{0.3333}}
  >{\raggedright\arraybackslash}p{(\linewidth - 4\tabcolsep) * \real{0.3333}}@{}}
\toprule\noalign{}
\begin{minipage}[b]{\linewidth}\raggedright
Abordagem
\end{minipage} & \begin{minipage}[b]{\linewidth}\raggedright
Tipo
\end{minipage} & \begin{minipage}[b]{\linewidth}\raggedright
Uso no ensaio
\end{minipage} \\
\midrule\noalign{}
\endhead
\bottomrule\noalign{}
\endlastfoot
Pérez et al.~(2006) & Local, paramétrica, mesma amostra MCMC &
\textbf{Sim} --- derivadas e MCSE \\
PSIS (Vehtari et al., 2024) & Reponderação estabilizada + diagnóstico &
\textbf{Sim} --- varreduras globais e \(\hat k\) \\
\emph{Power-scaling} (Kallioinen et al.) & Local, perturbação por
potências da priori/verossimilhança & \textbf{Fluxo} (priori e
verossimilhança separadas) \\
REPS (Ho, 2023) & Global, não paramétrica (bola de KL) & Não ---
referência para limitações das análises locais \\
\end{longtable}
}


    % Add a bibliography block to the postdoc
    
    
    
\end{document}
